\documentclass[11pt]{article}

\usepackage[a4paper,margin=1in]{geometry}
\usepackage{amsmath,amssymb,amsthm,mathtools}
\usepackage{booktabs,array,longtable}
\usepackage{enumitem}
\usepackage{microtype}
\usepackage{xcolor}
\usepackage{hyperref}
\usepackage[nameinlink,capitalise,noabbrev]{cleveref}

\hypersetup{
  colorlinks=true,
  linkcolor=blue!55!black,
  citecolor=blue!55!black,
  urlcolor=blue!55!black
}

\newtheorem{theorem}{Theorem}[section]
\newtheorem{proposition}[theorem]{Proposition}
\newtheorem{lemma}[theorem]{Lemma}
\newtheorem{corollary}[theorem]{Corollary}
\newtheorem{definition}[theorem]{Definition}
\newtheorem{remark}[theorem]{Remark}

\newcommand{\one}{\mathbf 1}
\newcommand{\R}{\mathbb R}
\newcommand{\ip}[2]{\left\langle #1,#2\right\rangle}
\newcommand{\norm}[1]{\left\lVert #1\right\rVert}
\newcommand{\Tr}{\operatorname{Tr}}
\newcommand{\pos}[1]{\left(#1\right)_{+}}
\newcommand{\cH}{\mathcal H}
\newcommand{\cQ}{\mathcal Q}
\newcommand{\cL}{\mathcal L}
\newcommand{\dd}{\,\mathrm d}
\newcommand{\E}{\mathbb E}
\newcommand{\GammaLaw}{\operatorname{Gamma}}
\newcommand{\BetaLaw}{\operatorname{Beta}}
\newcommand{\supp}{\operatorname{supp}}

\title{A Complete Proof of the Goodman--Style Lower Bound\\
       for Odd-Cycle Densities in Graphons}
\author{Consolidated working note for H.~Yang and collaborators}
\date{July 2026}


\begin{document}
\maketitle

\begin{abstract}
Let \(W:[0,1]^2\to[0,1]\) be a graphon of edge density \(p\), and let
\(m\ge3\) be odd.  We prove the conjectured Goodman--style inequality
\[
       t(C_m,W)\ge p^m-p(1-p)^{m-1}
\]
for every \(p\in[0,1]\).  The proof is divided at \(p=2/3\), or equivalently
at the complement density \(q=1/3\).

For \(p\ge2/3\), a two-sided Schur-complement identity cancels the unknown
odd moments of the compressed complement operator.  Expansion in complete
homogeneous symmetric polynomials is then diagonalized by a Dirichlet
average, converted to a beta integral, and smoothed by a Laplace transform
into expectations of a shifted gamma random variable.  The remaining
positivity follows from a uniform odd-to-even gamma moment inequality.

For \(1/2<p<2/3\), the Hilbert--Schmidt budget allows at most one compression
eigenvalue above \(q=1-p\).  Isolating this frontier eigenvalue and retaining
both its direct coupling and the coupling forced into the safe spectral
subspace leads to a one-dimensional Huber-type envelope.  Two canonical
dual witnesses settle its quadratic and linear branches.  The only finite
residue is a pair of exact univariate Bernstein certificates at \(m=9\).
The cycles \(C_3,C_5,C_7\) are also given separate elementary proofs.

A final section extracts the reusable mechanisms: one-step Jacobi
coefficient stripping, first-return and cyclic excursion expansions,
Dirichlet probabilistic polarization, beta--gamma/Laplace positivity
transfer, the Laguerre differential equation hidden in the gamma moment
lemma, and the one-dangerous-mode/Huber principle.  It also explains why the
same threshold \(q=1/3\) emerges independently from the two sides.
\end{abstract}

\tableofcontents

\section{The theorem, notation, and proof map}\label{sec:intro}

For a finite graph \(F\) and a graphon \(W\), write
\[
 t(F,W)=\int_{[0,1]^{V(F)}}
          \prod_{ij\in E(F)}W(x_i,x_j)
          \prod_{i\in V(F)}\dd x_i.
\]
In particular,
\[
 p=t(K_2,W)=\int_{[0,1]^2}W(x,y)\dd x\dd y
\]
is the edge density.  We prove the following theorem.

\begin{theorem}[Odd-cycle Goodman bound]\label{thm:main}
Let \(W\) be a graphon of edge density \(p\in[0,1]\).  For every odd integer
\(m\ge3\),
\begin{equation}\label{eq:main-goal}
       t(C_m,W)\ge p^m-p(1-p)^{m-1}.
\end{equation}
\end{theorem}

Put
\[
       U=1-W,\qquad q=t(K_2,U)=1-p.
\]
If \(p\le1/2\), then
\[
 p^m-pq^{m-1}=p\bigl(p^{m-1}-q^{m-1}\bigr)\le0,
\]
whereas \(t(C_m,W)\ge0\).  Thus only \(p>1/2\), or \(q<1/2\), needs proof.
The nontrivial range splits into
\[
 \underbrace{0\le q\le\frac13}_{p\ge2/3}
 \qquad\text{and}\qquad
 \underbrace{\frac13<q<\frac12}_{1/2<p<2/3}.
\]

\subsection{Proof map}

The argument has four modules.

\begin{enumerate}[label=\textbf{Module \arabic*.},leftmargin=3.4em]
\item \Cref{sec:short-cycles} gives elementary proofs for \(m=3,5,7\), at
      every edge density.
\item \Cref{sec:common-operator,sec:dense-region} proves all odd lengths when
      \(q\le1/3\).  The principal chain is
      \[
      \begin{gathered}
      \text{Schur cancellation}
      \longrightarrow \text{Dirichlet diagonalization}\\
      \longrightarrow \text{beta integral}
      \longrightarrow \text{gamma smoothing}.
      \end{gathered}
      \]
\item \Cref{sec:regionII-operator,sec:huber,sec:chart,sec:quadratic,sec:linear}
      proves all odd \(m\ge9\) when \(1/3<q<1/2\).  The principal chain is
      \[
      \begin{gathered}
      \text{one frontier eigenvalue}
      \longrightarrow \text{two coupling channels}\\
      \longrightarrow \text{Huber duality}
      \longrightarrow \text{two scalar witnesses}.
      \end{gathered}
      \]
\item \Cref{sec:transferable} isolates general principles and connects them
      to Jacobi matrices, orthogonal polynomials, continued fractions,
      lattice paths, Stein identities, and Laguerre polynomials.
\end{enumerate}

The final assembly is in \Cref{sec:assembly}.  No numerical optimization or
interval covering is used.  The two fixed Bernstein coefficient lists in
\Cref{app:bernstein} are exact rational certificates.

\subsection{Step-graphon reduction}\label{subsec:step-reduction}

\begin{lemma}[Finite-rank reduction]\label[lemma]{lem:step-reduction}
For each fixed odd \(m\ge3\), it is enough to prove \eqref{eq:main-goal} for
symmetric step graphons.
\end{lemma}

This is a standard graphon-approximation principle.  The interpretation of
step graphons as finite weighted graphs and the continuity of homomorphism
densities are developed in Lov\'asz~\cite[Sections~7.2 and~10.5, especially
Lemma~10.23]{Lovasz2012}; the dyadic conditional-expectation approximation
used below is also the elementary martingale approximation discussed in
\cite[Appendix~A.3.6]{Lovasz2012}.  We include the details because the present
argument needs the approximating graphons to have edge density exactly \(p\),
not merely a density converging to \(p\).

\begin{proof}
We spell out the approximation, including the meaning of the finite
sigma-algebras.  For \(n\ge1\), divide \([0,1]\) into the \(2^n\) dyadic
intervals
\[
 I_{n,a}=\left[\frac{a}{2^n},\frac{a+1}{2^n}\right),
 \qquad 0\le a<2^n,
\]
with the harmless convention that the last interval contains the endpoint
\(1\).  Let \(\mathcal A_n\) be the finite sigma-algebra consisting of all
unions of these intervals.  Thus the atoms of the product sigma-algebra
\(\mathcal A_n\otimes\mathcal A_n\) are precisely the dyadic rectangles
\(I_{n,a}\times I_{n,b}\).  The partitions refine as \(n\) increases, and
dyadic intervals generate the Borel sigma-algebra.

On each such rectangle, replace \(W\) by its average.  More explicitly, set
\begin{equation}\label{eq:dyadic-step-approximation}
 w^{(n)}_{ab}
   =2^{2n}\int_{I_{n,a}\times I_{n,b}}W(x,y)\dd x\dd y,
 \qquad
 W_n(x,y)=w^{(n)}_{ab}
 \quad\text{if }(x,y)\in I_{n,a}\times I_{n,b}.
\end{equation}
This cell-averaging operation is the standard conditional 
expectation operation with respect to the sigma algebra
$\mathcal A_n\otimes\mathcal A_n$, and is 
denoted by the notation
\(\E[\,\_\, \mid\mathcal A_n\otimes\mathcal A_n]\). So,
\(W_n = \E[W\mid\mathcal A_n\otimes\mathcal A_n] \). Since \(0\le W\le1\),
every cell average lies in \([0,1]\).  Moreover, symmetry of \(W\) gives
\(w^{(n)}_{ab}=w^{(n)}_{ba}\).  Hence \(W_n\) is a symmetric step graphon.

There is also a completely finite interpretation.  The matrix
\((w^{(n)}_{ab})_{0\le a,b<2^n}\) is the weighted adjacency matrix of a graph
on \(2^n\) equally weighted vertices (loops are allowed), and
\[
 t(C_m,W_n)
 =\frac{1}{2^{nm}}
   \sum_{a_1,\ldots,a_m=0}^{2^n-1}
   w^{(n)}_{a_1a_2}w^{(n)}_{a_2a_3}\cdots
   w^{(n)}_{a_ma_1}.
\]
The same finite structure can be seen directly at the operator level.  For
any function \(f\), if \(x\in I_{n,a}\), then
\[
 (T_{W_n}f)(x)
 =\int_0^1W_n(x,y)f(y)\dd y
 =\sum_{b=0}^{2^n-1}w^{(n)}_{ab}
       \int_{I_{n,b}}f(y)\dd y.
\]
Thus the output \(T_{W_n}f\) is constant on each interval \(I_{n,a}\), no
matter how complicated the input \(f\) is.  Every such output is therefore a
linear combination of the \(2^n\) indicator functions
\(\one_{I_{n,0}},\ldots,\one_{I_{n,2^n-1}}\).  The set of all possible
outputs is called the \emph{range} of the operator, and its dimension is the
\emph{rank}.  Since the range is contained in a space of dimension \(2^n\),
the rank is at most \(2^n\).  This explains the term ``finite-rank
reduction.''

Two facts remain to be checked: the edge density is unchanged, and the cycle
density converges.  For the first, summing the averages in
\eqref{eq:dyadic-step-approximation} over all cells gives
\[
 \int_{[0,1]^2}W_n(x,y)\dd x\dd y
 =\sum_{a,b}|I_{n,a}|\,|I_{n,b}|\,w^{(n)}_{ab}
 =\int_{[0,1]^2}W(x,y)\dd x\dd y=p.
\]
Thus \(t(K_2,W_n)=p\) for every \(n\), not merely in the limit.

Next we verify that \(W_n\to W\) in \(L^2([0,1]^2)\).  Dyadic rectangle
step functions are dense in this space.  Therefore, given \(\varepsilon>0\),
we may choose an integer \(N\) and a function \(S\), constant on every
rectangle \(I_{N,a}\times I_{N,b}\), such that
\(\norm{W-S}_2<\varepsilon\).  For \(n\ge N\), the function \(S\) is also
constant on every level-\(n\) rectangle.  Averaging is an \(L^2\)-contraction
(this is Jensen's inequality on each rectangle), so
\[
 \norm{W_n-S}_2
 =\norm{\E[W-S\mid\mathcal A_n\otimes\mathcal A_n]}_2
 \le \norm{W-S}_2.
\]
Consequently
\[
       \norm{W_n-W}_2
       \le \norm{W_n-S}_2+\norm{S-W}_2
       <2\varepsilon,
\]
which proves the asserted convergence.

Finally, cycle densities are continuous under this convergence.  Put \(x_{m+1}=x_1\), and for a fixed
\((x_1,\ldots,x_m)\), write
\[
       a_i=W_n(x_i,x_{i+1}),\qquad b_i=W(x_i,x_{i+1}).
\]
The telescoping identity
\[
 \prod_{i=1}^m a_i-\prod_{i=1}^m b_i
 =\sum_{j=1}^m
   \left(\prod_{i<j}a_i\right)(a_j-b_j)
   \left(\prod_{i>j}b_i\right)
\]
and the bounds \(0\le a_i,b_i\le1\) show, after integration, that
\begin{align*}
 \bigl|t(C_m,W_n)-t(C_m,W)\bigr|
 &\le \sum_{j=1}^m
       \int_{[0,1]^m}
       |W_n(x_j,x_{j+1})-W(x_j,x_{j+1})|
       \prod_{i=1}^m\dd x_i \\
 &=m\norm{W_n-W}_1
 \le m\norm{W_n-W}_2
 \longrightarrow0.
\end{align*}
Here the equality in the second line follows because all variables other
than \(x_j,x_{j+1}\) integrate out to \(1\).

Now suppose \eqref{eq:main-goal} has been proved for every symmetric step
graphon.  Applying it to \(W_n\), whose edge density is exactly \(p\), gives
\[
       t(C_m,W_n)\ge p^m-p(1-p)^{m-1}.
\]
Letting \(n\to\infty\) proves the same inequality for \(W\), as required.
\end{proof}

\begin{remark}[How to read the power-series calculations]
Several later arguments introduce an auxiliary variable \(z\) and use the
notation \([z^m]F(z)\).  The variable \(z\) is simply a bookkeeping device:
its exponent records the length of a walk or cycle.  If
\[
       F(z)=c_0+c_1z+c_2z^2+\cdots,
\]
then \([z^m]F(z)=c_m\).  Thus extracting \([z^m]\) means retaining exactly
the terms corresponding to length \(m\).

There are two equivalent rigorous ways to interpret the power series.  The
first is the \emph{formal} interpretation.  Here no numerical value is ever
assigned to \(z\), and convergence is not being asserted or needed.  We
manipulate coefficients according to the usual rules for generating
functions.  For example, for a finite matrix \(X\),
\[
       (I-zX)^{-1}=I+zX+z^2X^2+z^3X^3+\cdots
\]
is a valid formal identity: multiplying the right side by \(I-zX\) makes
all positive powers cancel coefficient by coefficient.  Similarly, if
\(Y(z)\) has zero constant term, then
\[
       -\log(1-Y(z))
       =\sum_{r\ge1}\frac{Y(z)^r}{r}
\]
is a formal identity.  For any fixed target coefficient \([z^m]\), only
finitely many powers on the right can contribute.  In this sense, the
subsequent calculations are finite algebraic calculations disguised as
power series.  Here ``formal'' is a technical algebraic term: two formal
series are equal when every coefficient is exactly equal.  

The second interpretation packages the same bookkeeping into ordinary
functions.  By \Cref{lem:step-reduction}, we may temporarily work with a
step graphon, whose operator is represented by a finite matrix \(X\).  At
\(z=0\), the matrix \(I-zX\) is just the identity matrix, so it is invertible.
It remains invertible when \(z\) is sufficiently close to \(0\).  In that
small neighbourhood, \((I-zX)^{-1}\), \(\det(I-zX)\), and the logarithm of
the determinant are honest ordinary functions with convergent Taylor
series.  Expanding those functions at \(z=0\) produces exactly the formal
coefficients described above.  The restriction to small \(|z|\) is used
only to justify these Taylor expansions; it imposes no additional condition
on the graphon.

Informally, the two interpretations are simply two ways of storing the same
list of coefficients.  The formal viewpoint manipulates the list directly;
the analytic viewpoint packages the list into a function near \(z=0\) and
then recovers it from the Taylor expansion.  Since the proof only asks for a
particular coefficient \([z^m]\), both viewpoints give the same exact finite
calculation.  We perform that calculation for the finite matrix associated
with a step graphon and then use \Cref{lem:step-reduction} to pass to the
original graphon.  Thus no determinant of an infinite-dimensional operator
is needed.
\end{remark}
\section{The elementary cases \texorpdfstring{\(m=3,5,7\)}{m=3,5,7}}\label{sec:short-cycles}

Throughout this section, \(W\) is any graphon, \(p=t(K_2,W)\),
\[
       U=1-W,\qquad q=t(K_2,U)=1-p,
\]
and \(T=T_U\) is the integral operator of \(U\).  When \(q>1/2\), the
right side of \eqref{eq:main-goal} is nonpositive, so the result is trivial.
We may therefore assume
\begin{equation}\label{eq:q-short}
       0\le q\le\frac12.
\end{equation}
The arguments below actually prove the three cases at every edge density.

\subsection{The triangle: Goodman's one-line argument}\label{subsec:C3}

The following is the graphon form of Goodman's classical counting
argument~\cite{Goodman1959}.

\begin{proposition}\label{prop:C3}
For every graphon \(W\) with edge density \(p=t(K_2,W)\),
\[
       t(C_3,W)\ge p^3-p(1-p)^2.
\]
\end{proposition}

\begin{proof}
The target polynomial is
\[
       p^3-p(1-p)^2=2p^2-p.
\]
Let \(d_W(x)=\int_0^1W(x,y)\dd y\).  Since
\(ab\ge a+b-1\) for \(a,b\in[0,1]\),
\[
 \int_0^1W(x,z)W(y,z)\dd z\ge d_W(x)+d_W(y)-1.
\]
Multiplying by \(W(x,y)\ge0\), integrating in \((x,y)\), and applying
Jensen's inequality gives
\[
\begin{aligned}
 t(C_3,W)
 &\ge \int W(x,y)\bigl(d_W(x)+d_W(y)-1\bigr)\dd x\dd y\\
 &=2\int_0^1d_W(x)^2\dd x-p
 \ge2p^2-p,
\end{aligned}
\]
as desired.
\end{proof}

\subsection{Path moments in the complement}\label{subsec:path-moments}

Write
\begin{equation}\label{eq:T1-decomp}
       T\one=q\one+g,
       \qquad g\perp\one.
\end{equation}
This decomposition has the explicit pointwise form
\[
 g(x)=\int_0^1U(x,y)\dd y-q
     =\bigl(1-d_W(x)\bigr)-(1-p)
     =p-d_W(x).
\]
Thus \(g\) is the negative of the deviation of the \(W\)-degree from its
mean \(p\).
For \(j\ge1\), let \(P_j\) denote the path with \(j\) edges and put
\[
       x_j=t(P_j,U)=\ip{\one}{T^j\one}.
\]
Thus \(x_1=q\).  The following first moments will be used for \(C_5\).

\begin{lemma}\label{lem:path-to-4}
With \eqref{eq:T1-decomp},
\[
\begin{aligned}
 x_2&=q^2+\norm g^2,\\
 x_3&=q^3+2q\norm g^2+\ip g{Tg},\\
 x_4&=q^4+3q^2\norm g^2+2q\ip g{Tg}+\norm{Tg}^2.
\end{aligned}
\]
\end{lemma}

\begin{proof}
The first identity is \(x_2=\norm{T\one}^2\).  Next,
\(T^2\one=q^2\one+qg+Tg\), while
\[
       \ip{\one}{Tg}=\ip{T\one}{g}=\norm g^2.
\]
Expanding \(x_3=\ip{T\one}{T^2\one}\) and
\(x_4=\norm{T^2\one}^2\) gives the displayed formulas.
\end{proof}

For \(C_7\), let \(H_0=\one^\perp\) and
\[
       A=P_{H_0}TP_{H_0}.
\]
Concretely, \(H_0\) is the space of mean-zero functions, and its orthogonal
projection is the centering operation
\[
       (P_{H_0}f)(x)=f(x)-\int_0^1f(y)\dd y.
\]
Thus \(P_{H_0}\) removes the constant part of \(f\), while \(A\) centers the
input, applies \(T\), and then centers the output.  In particular, \(A\) acts
entirely within the mean-zero subspace.
Define
\begin{equation}\label{eq:sj-def}
       s_j=\ip g{A^jg}\qquad(j\ge0).
\end{equation}
The spectral theorem supplies a finite positive measure \(\mu\) on the real
line such that
\[
       s_j=\int\lambda^j\dd\mu(\lambda),
       \qquad s_0=\norm g^2=\mu(\R).
\]

A useful way to verify all path formulas is to write
\[
       T^n\one=a_n\one+h_n,\qquad h_n\perp\one.
\]
Then \(a_n=x_n\), and the block form of \(T\) gives the recursion
\begin{equation}\label{eq:path-recursion}
       a_{n+1}=qa_n+\ip g{h_n},
       \qquad h_{n+1}=a_ng+Ah_n,
       \qquad (a_0,h_0)=(1,0).
\end{equation}
For clarity, iterating the second recursion first gives
\[
\begin{aligned}
 h_1={}&g,\\
 h_2={}&qg+Ag,\\
 h_3={}&a_2g+qAg+A^2g,\\
 h_4={}&a_3g+a_2Ag+qA^2g+A^3g,\\
 h_5={}&a_4g+a_3Ag+a_2A^2g+qA^3g+A^4g,\\
 h_6={}&a_5g+a_4Ag+a_3A^2g+a_2A^3g
          +qA^4g+A^5g.
\end{aligned}
\]
More generally,
\[
       h_n=\sum_{j=0}^{n-1}a_{n-1-j}A^jg.
\]
Taking the inner product with \(g\), and recalling that
\(s_j=\ip g{A^jg}\), turns the first recursion in
\eqref{eq:path-recursion} into the scalar formula
\begin{equation}\label{eq:path-scalar-recursion}
       a_{n+1}=qa_n+\sum_{j=0}^{n-1}a_{n-1-j}s_j
       \qquad(n\ge1).
\end{equation}
For example, since \(a_0=1\) and \(a_1=q\),
\[
\begin{aligned}
 a_2&=qa_1+a_0s_0=q^2+s_0,\\
 a_3&=qa_2+a_1s_0+a_0s_1=q^3+2qs_0+s_1,\\
 a_4&=qa_3+a_2s_0+a_1s_1+a_0s_2
      =q^4+3q^2s_0+2qs_1+s_0^2+s_2.
\end{aligned}
\]
Continuing the same calculation gives
\begin{equation}\label{eq:path-2-6}
\begin{aligned}
 x_2={}&a_2 = q^2+s_0,\\
 x_3={}&a_3 = q^3+2qs_0+s_1,\\
 x_4={}&a_4 = q^4+3q^2s_0+2qs_1+s_0^2+s_2,\\
 x_5={}&a_5 = q^5+4q^3s_0+3q^2s_1+3qs_0^2+2qs_2+2s_0s_1+s_3,\\
 x_6={}&a_6 = q^6+5q^4s_0+4q^3s_1+6q^2s_0^2+3q^2s_2
       +6qs_0s_1+2qs_3\\
     &\qquad{}+s_0^3+2s_0s_2+s_1^2+s_4.
\end{aligned}
\end{equation}

\subsection{The pentagon: one completion of a square}\label{subsec:C5}

\begin{proposition}\label{prop:C5}
For every graphon \(W\) with edge density \(p=t(K_2,W)\),
\[
       t(C_5,W)\ge p^5-p(1-p)^4.
\]
\end{proposition}

\begin{proof}
Assume \eqref{eq:q-short}.  Label the integration variables cyclically by
\(z_1,\ldots,z_5\), put \(z_6=z_1\), and abbreviate
\[
       U_i=U(z_i,z_{i+1})\qquad(1\le i\le5).
\]
Since \(W=1-U\), the integrand defining \(t(C_5,W)\) is
\(\prod_{i=1}^5(1-U_i)\).  Ordinary inclusion--exclusion gives the exact
pointwise expansion
\[
\begin{aligned}
 \prod_{i=1}^5(1-U_i)
 ={}&1-\sum_{i=1}^5U_i
     +\sum_{1\le i<j\le5}U_iU_j
     -\sum_{1\le i<j<k\le5}U_iU_jU_k\\
    &+\sum_{i=1}^5\prod_{j\ne i}U_j
     -\prod_{i=1}^5U_i.
\end{aligned}
\]
We now integrate this identity over \((z_1,\ldots,z_5)\).  Each single-edge
term has integral \(q=x_1\), so the first sum contributes \(5q\).  Among the
\(\binom52=10\) pairs of edges, five pairs are adjacent around the cycle;
each forms a two-edge path and hence has integral \(x_2\).  The other five
pairs are disjoint edges.  Their variables separate, so each has integral
\(x_1^2=q^2\).  Therefore
\[
       \int\sum_{i<j}U_iU_j=5x_2+5q^2.
\]

Similarly, among the \(\binom53=10\) triples, five consist of three
consecutive edges and form a path with three edges, contributing \(x_3\)
each.  Each of the other five consists of a two-edge path together with one
disjoint edge.  Homomorphism density is multiplicative over disjoint unions,
so each of these contributes \(x_2x_1=qx_2\).  Thus
\[
       \int\sum_{i<j<k}U_iU_jU_k=5x_3+5qx_2.
\]
There are five ways to select four edges; deleting any one edge from the
cycle leaves a four-edge path, so every corresponding integral is \(x_4\).
Finally, selecting all five edges gives \(c_5:=t(C_5,U)\).  Combining these
contributions with their alternating signs yields
\[
\begin{aligned}
 t(C_5,1-U)
  ={}&1-5q+5x_2+5q^2-5x_3-5qx_2+5x_4-c_5.
\end{aligned}
\]
Since \(0\le U\le1\), deleting one edge from the cycle integrand gives
\(c_5\le x_4\).  Therefore
\begin{equation}\label{eq:C5-IE}
 t(C_5,W)\ge1-5q+5q^2+5(1-q)x_2-5x_3+4x_4.
\end{equation}
The target is
\[
 (1-q)^5-(1-q)q^4=1-5q+10q^2-10q^3+4q^4.
\]
Subtracting it from \eqref{eq:C5-IE} and using
\Cref{lem:path-to-4} gives
\begin{equation}\label{eq:Phi5}
\begin{aligned}
 \Phi_5={}&4\norm{Tg}^2+(8q-5)\ip g{Tg}
        +(12q^2-15q+5)\norm g^2\\
 ={}&4\left\|Tg+\frac{8q-5}{8}g\right\|^2
   +\left(8q^2-10q+\frac{55}{16}\right)\norm g^2.
\end{aligned}
\end{equation}
The residual quadratic has discriminant
\[
       100-4\cdot8\cdot\frac{55}{16}=-10<0
\]
and positive leading coefficient.  Thus \(\Phi_5\ge0\), proving the claim.
\end{proof}

\subsection{The heptagon: a univariate sum of squares}\label{subsec:C7}

The final certificate below is an explicit one-variable sum of squares; for
general background on positivity of polynomials on an interval, see
Powers--Reznick~\cite{PowersReznick2000}.

\begin{proposition}\label{prop:C7}
For every graphon \(W\) with edge density \(p=t(K_2,W)\),
\[
       t(C_7,W)\ge p^7-p(1-p)^6.
\]
\end{proposition}

\begin{proof}
Assume \eqref{eq:q-short}.  As in the proof for \(C_5\), label the integration
variables cyclically by \(z_1,\ldots,z_7\), put \(z_8=z_1\), and write
\[
       U_i=U(z_i,z_{i+1})\qquad(1\le i\le7).
\]
Expanding the seven factors \(1-U_i\) gives the exact inclusion--exclusion
identity
\begin{equation}\label{eq:C7-subset-expansion}
 t(C_7,W)
 =\sum_{S\subseteq\{1,\ldots,7\}}(-1)^{|S|}
   \int_{[0,1]^7}\prod_{i\in S}U_i\prod_{j=1}^7\dd z_j.
\end{equation}
For a proper subset \(S\), the selected edges break into disjoint runs of
consecutive edges around the cycle.  If their lengths are
\(r_1,\ldots,r_\ell\), the selected subgraph is the disjoint union of paths
\(P_{r_1}\sqcup\cdots\sqcup P_{r_\ell}\).  Its variables separate between
components, so its contribution to \eqref{eq:C7-subset-expansion} is
\[
       x_{r_1}\cdots x_{r_\ell}.
\]
In particular, a run of length one contributes \(x_1=q\).  The proper edge
subsets of the labelled \(7\)-cycle have the following run-length types.
\[
\begin{array}{c|l}
 k&\text{types inside parenthesis and their multiplicities outside them for }k\text{ selected edges}\\
\hline
0&1\\
1&7(1)\\
2&7(2)+14(1+1)\\
3&7(3)+21(2+1)+7(1+1+1)\\
4&7(4)+14(3+1)+7(2+2)+7(2+1+1)\\
5&7(5)+7(4+1)+7(3+2)\\
6&7(6).
\end{array}
\]
Here, for example, \(14(3+1)\) denotes 
the fourteen different ways (i.e., multiplicity 14) to select a three-edge path together with a
disjoint edge.  The multiplicities can be checked directly around the
cycle.  For two selected edges, seven pairs are adjacent and the remaining
\(21-7=14\) are disjoint.  For three selected edges, there are seven
rotations of a three-edge run and seven rotations of the alternating pattern
with three isolated edges; the remaining \(35-7-7=21\) have type \(2+1\).
For four selected edges, there are seven four-edge runs, fourteen choices of
type \(3+1\), and seven choices of type \(2+2\); the remaining seven have
type \(2+1+1\).  Finally, for five selected edges, the two omitted edges are
either adjacent, giving seven five-edge runs, or nonadjacent, giving seven
configurations of each type \(4+1\) and \(3+2\).  Thus the multiplicities in
each row sum, as they should, to \(\binom7k\).

Substituting the path-density product associated with each entry of the
table into \eqref{eq:C7-subset-expansion} gives
\begin{align*}
 t(C_7,W)={}&1-7q
 +(7x_2+14q^2)
 -(7x_3+21qx_2+7q^3)\\
 &+(7x_4+14qx_3+7x_2^2+7q^2x_2)\\
 &-(7x_5+7qx_4+7x_3x_2)
 +7x_6-c_7,
\end{align*}
where \(c_7=t(C_7,U)\) is the contribution obtained by selecting all seven
edges.  Because \(0\le U_i\le1\), deleting one factor from that integrand
can only increase it; the remaining six factors form a six-edge path.
Hence
\[
       c_7=\int\prod_{i=1}^7U_i
       \le\int\prod_{i=1}^6U_i=x_6.
\]
Replacing \(-c_7\) by \(-x_6\) in the preceding identity and collecting the
coefficients of the same path densities gives
\begin{equation}\label{eq:C7-IE}
\begin{aligned}
 t(C_7,W)\ge{}&6x_6-7x_5+7(1-q)x_4+7(2q-1)x_3\\
 &+7(q^2-3q+1)x_2+7x_2^2-7x_2x_3\\
 &+1-7q+14q^2-7q^3.
\end{aligned}
\end{equation}
Since
\[
 (1-q)^7-(1-q)q^6
 =1-7q+21q^2-35q^3+35q^4-21q^5+6q^6,
\]
substitution of \eqref{eq:path-2-6} reduces the desired inequality to
\begin{equation}\label{eq:Phi7-expanded}
\begin{aligned}
 \Phi_7={}&6s_4+(12q-7)s_3+(18q^2-21q+7)s_2\\
 &+(24q^3-42q^2+28q-7)s_1\\
 &+(30q^4-70q^3+70q^2-35q+7)s_0\\
 &+12s_0s_2+(36q-21)s_0s_1
 +(36q^2-42q+14)s_0^2+6s_0^3+6s_1^2.
\end{aligned}
\end{equation}
Put \(r=s_0=\mu(\R)\).  Grouping the terms integrated against \(\mu\),
\begin{equation}\label{eq:Phi7-integral}
       \Phi_7=6s_1^2+\int R_{q,r}(\lambda)\dd\mu(\lambda),
\end{equation}
where
\[
       R_{q,r}(\lambda)=P_q(\lambda)+rB_q(\lambda)+6r^2,
\]
\begin{align*}
 P_q(\lambda)={}&6\lambda^4+(12q-7)\lambda^3
 +(18q^2-21q+7)\lambda^2\\
 &+(24q^3-42q^2+28q-7)\lambda
 +30q^4-70q^3+70q^2-35q+7,\\
 B_q(\lambda)={}&12\lambda^2+(36q-21)\lambda+36q^2-42q+14.
\end{align*}
We show that both \(P_q\) and \(B_q\) are nonnegative.

The minimum of \(B_q\) as a quadratic in \(\lambda\) occurs at
\(\lambda_*=(7-12q)/8\), and
\[
 B_q(\lambda_*)=\frac{144q^2-168q+77}{16}\ge\frac{29}{16},
\]
because
\[
 144q^2-168q+48=144\left(q-\frac12\right)
                         \left(q-\frac23\right)\ge0
\]
on \([0,1/2]\).

For \(P_q\), set
\begin{align*}
 D(q)&=288q^2-336q+119,\\
 C(q)&=24q^3-42q^2+28q-7,\\
 N(q)&=5184q^6-18144q^5+28602q^4-25802q^3\\
     &\qquad+13874q^2-4165q+539.
\end{align*}
Direct expansion gives the exact identity
\begin{equation}\label{eq:Pq-SOS}
\begin{aligned}
 P_q(\lambda)
 ={}&6\left(\lambda^2+\left(q-\frac7{12}\right)\lambda\right)^2\\
 &+\frac{D(q)}{24}
   \left(\lambda+\frac{12C(q)}{D(q)}\right)^2
 +\frac{N(q)}{D(q)}.
\end{aligned}
\end{equation}
Writing \(y=1/2-q\ge0\),
\[
       D(q)=288y^2+48y+23>0
\]
and
\[
\begin{aligned}
 8N(q)={}&41472y^6+20736y^5+21456y^4+7984y^3\\
         &+2032y^2+316y+11>0.
\end{aligned}
\]
Thus \(P_q(\lambda)\ge0\) for every real \(\lambda\).  Since
\(r\ge0\), \(R_{q,r}(\lambda)\ge0\), and \eqref{eq:Phi7-integral} yields
\(\Phi_7\ge0\).
\end{proof}

\begin{remark}
The \(C_5\) and \(C_7\) arguments use no positivity of the kernel operator
\(T_U\); a pointwise nonnegative kernel need not define a positive
semidefinite operator.  Only self-adjointness, the positive spectral measure
associated with a vector, and the displayed algebraic identities are used.
\end{remark}


\section{The common complement-operator framework}\label{sec:common-operator}

By \Cref{lem:step-reduction}, determinant arguments may be carried out for a
step graphon.  Decompose
\[
       L^2([0,1])=\R\one\oplus\one^\perp
\]
and write the complement operator as
\begin{equation}\label{eq:block-U}
       T_U=
       \begin{pmatrix}
          q&g^*\\
          g&A
       \end{pmatrix},
       \qquad
       g=T_U\one-q\one,
       \qquad
       A=P_{\one^\perp}T_UP_{\one^\perp}.
\end{equation}
The vector \(\one\) has norm one because the underlying probability space
has total mass one.

\begin{lemma}[Square budget]\label{lem:HS-budget}
One has
\begin{equation}\label{eq:HS-budget}
       \Tr(A^2)+2\norm g^2\le pq.
\end{equation}
\end{lemma}

\begin{proof}
Since \(0\le U\le1\),
\[
 \Tr(T_U^2)=\int_{[0,1]^2}U(x,y)^2\dd x\dd y
 \le\int_{[0,1]^2}U(x,y)\dd x\dd y=q.
\]
On the other hand, \eqref{eq:block-U} gives
\[
       \Tr(T_U^2)=q^2+2\norm g^2+\Tr(A^2).
\]
Subtracting \(q^2\) proves the claim.
\end{proof}

\begin{lemma}[Half-norm lemma]\label{lem:half-norm}
The spectrum of \(A\) lies in \([-1/2,1/2]\); equivalently,
\begin{equation}\label{eq:A-half}
       \norm A_{\mathrm{op}}\le\frac12.
\end{equation}
\end{lemma}

\begin{proof}
Let \(f\perp\one\) and \(\norm f=1\), and write \(f=f_+-f_-\).  Since
\(\int f=0\),
\[
       \int f_+=\int f_-=:a,
       \qquad 2a=\lVert f\rVert_1\le\lVert f\rVert_2=1.
\]
Using \(0\le U\le1\),
\[
 \ip f{T_Uf}
 \le\left(\int f_+\right)^2+\left(\int f_-\right)^2
 =2a^2\le\frac12,
\]
and similarly
\[
 \ip f{T_Uf}
 \ge-2\left(\int f_+\right)\left(\int f_-\right)
 =-2a^2\ge-\frac12.
\]
The variational characterization of the norm of a self-adjoint operator
proves the result.
\end{proof}

\begin{remark}[No positivity-of-operator assumption]
The pointwise inequality \(U\ge0\) does \emph{not} imply that \(T_U\) or its
compression \(A\) is positive semidefinite.  The proof above controls the
quadratic form directly and is the only spectral support statement needed.
\end{remark}

Let \(E_A\) be the spectral resolution of \(A\), and define the finite
positive measure
\begin{equation}\label{eq:mu-def}
       \mu(B)=\ip g{E_A(B)g}.
\end{equation}
Then \(\mu\) is supported on \([-1/2,1/2]\), and
\begin{equation}\label{eq:resolvents-common}
 \ip g{(I-zA)^{-1}g}
 =\int_{-1/2}^{1/2}\frac{\dd\mu(\lambda)}{1-\lambda z}.
\end{equation}
This positive spectral measure is the bridge from the graphon operator to
one-variable moment inequalities.

\section{The dense region \texorpdfstring{\(p\ge2/3\)}{p>=2/3}: Dirichlet and gamma smoothing}
\label{sec:dense-region}

Throughout this section,
\begin{equation}\label{eq:dense-q}
       0\le q\le\frac13,
       \qquad p=1-q,
\end{equation}
and \(m\ge3\) is odd.  We prove \eqref{eq:main-goal} for all such \(m\).
The proof is first presented as a sequence of exact reductions; a conceptual
interpretation is given in \Cref{sec:transferable}.

\subsection{A two-sided spectral shift and cancellation of the tail}

Define
\begin{align}
 R_-(z)&=\ip g{(I+zA)^{-1}g}
       =\int\frac{\dd\mu(\lambda)}{1+\lambda z},\label{eq:dense-Rminus}\\
 R_+(z)&=\ip g{(I-zA)^{-1}g}
       =\int\frac{\dd\mu(\lambda)}{1-\lambda z},\label{eq:dense-Rplus}
\end{align}
and
\begin{align}
 L_-(z)&=-\log\left(1-\frac{z^2R_-(z)}{1-pz}\right),\label{eq:dense-Lminus}\\
 L_+(z)&=-\log\left(1-\frac{z^2R_+(z)}{1-qz}\right).\label{eq:dense-Lplus}
\end{align}
For odd \(m\), put
\begin{equation}\label{eq:dense-Sm}
       S_m=m[z^m]\bigl(L_-(z)+L_+(z)\bigr).
\end{equation}

\begin{lemma}[Two-sided odd cancellation]\label{lem:dense-cancellation}
For every odd \(m\ge3\),
\begin{equation}\label{eq:dense-cycle-sum}
       t(C_m,W)+t(C_m,U)=p^m+q^m+S_m.
\end{equation}
Consequently,
\begin{equation}\label{eq:dense-defect}
 t(C_m,W)-\bigl(p^m-pq^{m-1}\bigr)
 =q^{m-1}+S_m-t(C_m,U).
\end{equation}
\end{lemma}

\begin{proof}
The operator of \(W=1-U\) is
\[
 T_W=
 \begin{pmatrix}p&-g^*\\-g&-A\end{pmatrix}.
\]
Conjugating by the sign flip on \(\one^\perp\), it is unitarily equivalent
to
\[
       M=\begin{pmatrix}p&g^*\\g&-A\end{pmatrix}.
\]
Schur complementation gives
\begin{align*}
 \det(I-zM)
 &=\det(I+zA)(1-pz)
   \left(1-\frac{z^2R_-(z)}{1-pz}\right),\\
 \det(I-zT_U)
 &=\det(I-zA)(1-qz)
   \left(1-\frac{z^2R_+(z)}{1-qz}\right).
\end{align*}
Using
\[
       -\log\det(I-zX)
       =\sum_{j\ge1}\frac{\Tr(X^j)}{j}z^j,
\]
extract the coefficient of \(z^m\) and add the two formulas.  Because \(m\)
is odd, the contributions of \(\det(I+zA)\) and \(\det(I-zA)\) cancel.
The scalar factors contribute \(p^m+q^m\), and the remaining factors
contribute \(S_m\).  This proves \eqref{eq:dense-cycle-sum}.  Finally,
\(q^m+pq^{m-1}=q^{m-1}\), giving \eqref{eq:dense-defect}.
\end{proof}

The cancellation is the decisive reason for using \emph{both} \(W\) and its
complement in this region: all odd moments \(\Tr(A^m)\) of the inaccessible
compressed tail disappear.  For background on determinant/trace expansions
and finite-rank perturbations, see Simon~\cite{Simon2005}.

\subsection{Replacing a complement cycle by a path}

Let \(P_j\) denote the path with \(j\) edges and set
\begin{equation}\label{eq:dense-xj}
       x_j=t(P_j,U)=\ip\one{T_U^j\one}.
\end{equation}
Deleting one factor from the integrand of \(t(C_m,U)\), and using \(0\le U\le1\), gives
\begin{equation}\label{eq:dense-cycle-path}
       t(C_m,U)\le x_{m-1}.
\end{equation}
Block inversion in \eqref{eq:block-U} gives the moment generating function
\begin{equation}\label{eq:dense-path-gf}
 \sum_{j\ge0}x_jz^j
 =\ip\one{(I-zT_U)^{-1}\one}
 =\frac1{1-qz-z^2R_+(z)}.
\end{equation}
Thus it is enough to prove
\begin{equation}\label{eq:dense-Phi}
       \Phi_m:=q^{m-1}+S_m-x_{m-1}\ge0.
\end{equation}

\subsection{Expansion in complete homogeneous polynomials}

For variables \(y_1,\ldots,y_k\), define
\[
 h_d(y_1,\ldots,y_k)
 =\sum_{\substack{\alpha_1,\ldots,\alpha_k\ge0\\
                   \alpha_1+\cdots+\alpha_k=d}}
        y_1^{\alpha_1}\cdots y_k^{\alpha_k},
 \qquad h_0=1.
\]
Equivalently,
\begin{equation}\label{eq:dense-h-gf}
       \prod_{i=1}^k\frac1{1-y_i z}
       =\sum_{d\ge0}h_d(y_1,\ldots,y_k)z^d.
\end{equation}
The notation \(a^{\times r}\) means that \(a\) is repeated \(r\) times.

\begin{proposition}[Exact coefficient expansion]\label{prop:dense-expansion}
For odd \(m\ge3\),
\begin{equation}\label{eq:dense-expansion}
 \Phi_m=\sum_{r=1}^{(m-1)/2}
 \int\!\cdots\!\int
 P_{m,r}(q;\lambda_1,\ldots,\lambda_r)
 \prod_{i=1}^r\dd\mu(\lambda_i),
\end{equation}
where \(n=m-2r\) and
\begin{align}
 P_{m,r}(q;\lambda_1,\ldots,\lambda_r)
 ={}&\frac mr\Bigl[
 h_n(p^{\times r},-\lambda_1,\ldots,-\lambda_r)
 +h_n(q^{\times r},\lambda_1,\ldots,\lambda_r)
 \Bigr]\notag\\
 &-h_{n-1}(q^{\times(r+1)},\lambda_1,\ldots,\lambda_r).
 \label{eq:dense-Pmr}
\end{align}
\end{proposition}

\begin{proof}
Put
\[
       Y_-(z)=\frac{z^2R_-(z)}{1-pz},
       \qquad
       Y_+(z)=\frac{z^2R_+(z)}{1-qz}.
\]
From \(-\log(1-Y)=\sum_{r\ge1}Y^r/r\),
\[
 m[z^m]L_-(z)
 =\sum_{r\ge1}\frac mr
   [z^{m-2r}]\frac{R_-(z)^r}{(1-pz)^r}.
\]
The spectral representation gives
\[
 R_-(z)^r
 =\int\!\cdots\!\int
   \prod_{i=1}^r\frac1{1+\lambda_i z}
   \prod_{i=1}^r\dd\mu(\lambda_i).
\]
By \eqref{eq:dense-h-gf}, the \(r\)-th contribution is the integral of
\[
       \frac mr h_n(p^{\times r},-\lambda_1,\ldots,-\lambda_r).
\]
The same argument for \(L_+\) produces the second term in brackets in
\eqref{eq:dense-Pmr}.

On the other hand, \eqref{eq:dense-path-gf} expands as
\[
 \frac1{1-qz-z^2R_+(z)}
 =\sum_{r\ge0}\frac{(z^2R_+(z))^r}{(1-qz)^{r+1}}.
\]
The term \(r=0\) contributes \(q^{m-1}\) to \(x_{m-1}\), cancelling the
first term in \eqref{eq:dense-Phi}.  For \(r\ge1\), the coefficient of
\(z^{m-1}\) is the integral of
\[
       h_{n-1}(q^{\times(r+1)},\lambda_1,\ldots,\lambda_r).
\]
Combining the contributions proves the proposition.
\end{proof}

\subsection{Dirichlet diagonalization: a probabilistic polarization}

Define the diagonal polynomial
\begin{equation}\label{eq:dense-Ptilde}
       \widetilde P_{m,r}(q,\ell)
       =P_{m,r}(q;\ell,\ldots,\ell).
\end{equation}

\begin{lemma}[Dirichlet diagonalization]\label{lem:dense-dirichlet}
Let \((\Theta_1,\ldots,\Theta_r)\) be uniformly distributed on the simplex
\[
       \Theta_i\ge0,\qquad \sum_{i=1}^r\Theta_i=1.
\]
Then
\begin{equation}\label{eq:dense-dirichlet-identity}
 P_{m,r}(q;\lambda_1,\ldots,\lambda_r)
 =\E\,\widetilde P_{m,r}
       \left(q,\sum_{i=1}^r\Theta_i\lambda_i\right).
\end{equation}
Consequently,
\begin{equation}\label{eq:dense-diagonal-min}
 \min_{\lambda_i\in[-1/2,1/2]}
 P_{m,r}(q;\lambda_1,\ldots,\lambda_r)
 =\min_{\ell\in[-1/2,1/2]}\widetilde P_{m,r}(q,\ell).
\end{equation}
\end{lemma}

\begin{proof}
For every \(j\ge0\), the Dirichlet moment formula gives
\begin{equation}\label{eq:dense-dirichlet-moment}
 \E\left(\sum_{i=1}^r\Theta_i\lambda_i\right)^j
 =\frac{h_j(\lambda_1,\ldots,\lambda_r)}
        {\binom{j+r-1}{r-1}}.
\end{equation}
Indeed, expand the power and use
\[
 \E\prod_{i=1}^r\Theta_i^{\alpha_i}
 =\frac{(r-1)!\prod_i\alpha_i!}{(j+r-1)!},
 \qquad \sum_i\alpha_i=j.
\]
Also,
\begin{equation}\label{eq:dense-h-split}
 h_d(a^{\times k},b_1,\ldots,b_r)
 =\sum_{j=0}^d
   \binom{d-j+k-1}{k-1}a^{d-j}h_j(b_1,\ldots,b_r).
\end{equation}
Apply \eqref{eq:dense-h-split} to each term in \eqref{eq:dense-Pmr} and then
use \eqref{eq:dense-dirichlet-moment}; this gives
\eqref{eq:dense-dirichlet-identity}.  Every convex combination
\(\sum_i\Theta_i\lambda_i\) lies in \([-1/2,1/2]\), so the expectation is
at least the diagonal minimum.  Conversely, taking
\(\lambda_1=\cdots=\lambda_r=\ell\) realizes every diagonal value.
\end{proof}

Thus \eqref{eq:dense-expansion} will be nonnegative once we prove
\begin{equation}\label{eq:dense-diagonal-target}
       \widetilde P_{m,r}(q,\ell)\ge0
\end{equation}
for
\[
 0\le q\le\frac13,
 \quad 1\le r\le\frac{m-1}{2},
 \quad -\frac12\le\ell\le\frac12.
\]
This is the first major dimension reduction: an \(r\)-variable spectral
polynomial is replaced exactly, not approximately, by one scalar variable.

\subsection{A beta representation}

Fix \(r\ge1\) and put
\begin{equation}\label{eq:dense-n}
       n=m-2r.
\end{equation}
Then \(n\ge1\) is odd.  Let \(\Xi\sim\BetaLaw(r,r)\), and for \(x\in[0,1]\)
set
\begin{equation}\label{eq:dense-VW}
       V_x=qx+\ell(1-x),
       \qquad
       W_x=px-\ell(1-x).
\end{equation}
Thus \(V_x+W_x=x\).  Define
\begin{equation}\label{eq:dense-rho}
       \rho_{n,m}(u)
       =\frac mn\bigl(u^n+(1-u)^n\bigr)-u^{n-1}.
\end{equation}

\begin{lemma}[Beta formula]\label{lem:dense-beta}
One has
\begin{align}
 \widetilde P_{m,r}(q,\ell)
 ={}&\binom{n+2r-1}{2r-1}\frac nr
 \left[
   \frac mn\E\bigl(V_\Xi^n+W_\Xi^n\bigr)
   -\E\bigl(\Xi V_\Xi^{n-1}\bigr)
 \right]\label{eq:dense-beta-first}\\
 ={}&\binom{n+2r-1}{2r-1}\frac nr
 \E\left[
   \Xi^n\rho_{n,m}\left(q+\ell\frac{1-\Xi}{\Xi}\right)
 \right].\label{eq:dense-beta-second}
\end{align}
If \(\ell>0\), then
\begin{equation}\label{eq:dense-beta-integral}
 \widetilde P_{m,r}(q,\ell)
 =\frac{\Gamma(m)}{r\Gamma(n)\Gamma(r)^2}\,
  \ell^{n+r}
  \int_0^\infty
  \frac{s^{r-1}}{(\ell+s)^m}\rho_{n,m}(q+s)\dd s.
\end{equation}
\end{lemma}

\begin{proof}
For arbitrary \(a,b\), comparison of coefficients with the beta moments
shows
\begin{equation}\label{eq:dense-beta-h}
 h_n(a^{\times r},b^{\times r})
 =\binom{n+2r-1}{2r-1}
  \E\bigl(a\Xi+b(1-\Xi)\bigr)^n.
\end{equation}
Differentiating with respect to the common variable \(a\), and noting that
the derivative of the left side is
\(r h_{n-1}(a^{\times(r+1)},b^{\times r})\), gives
\begin{equation}\label{eq:dense-beta-h-deriv}
 h_{n-1}(a^{\times(r+1)},b^{\times r})
 =\binom{n+2r-1}{2r-1}\frac nr
  \E\left[\Xi\bigl(a\Xi+b(1-\Xi)\bigr)^{n-1}\right].
\end{equation}
Apply these formulas to the three terms of \eqref{eq:dense-Pmr}; the two
linear forms are precisely \(V_\Xi\) and \(W_\Xi\).  This proves
\eqref{eq:dense-beta-first}.  If
\[
       u=q+\ell\frac{1-x}{x},
\]
then \(V_x=xu\) and \(W_x=x(1-u)\), giving
\eqref{eq:dense-beta-second}.

For \(\ell>0\), use the beta density
\[
       \frac{\Gamma(2r)}{\Gamma(r)^2}
       x^{r-1}(1-x)^{r-1}\dd x
\]
and make the projective substitution
\[
       s=\ell\frac{1-x}{x},
       \qquad x=\frac\ell{\ell+s}.
\]
Then
\[
 x^{n+r-1}(1-x)^{r-1}\dd x
 =-\ell^{n+r}\frac{s^{r-1}}{(\ell+s)^m}\dd s.
\]
Substitution into \eqref{eq:dense-beta-second}, together with
\[
 \binom{m-1}{2r-1}\frac nr
 \frac{\Gamma(2r)}{\Gamma(r)^2}
 =\frac{\Gamma(m)}{r\Gamma(n)\Gamma(r)^2},
\]
proves \eqref{eq:dense-beta-integral}.
\end{proof}

\subsection{The nonpositive diagonal is pointwise safe}

\begin{lemma}\label{lem:dense-ell-negative}
If \(0\le q\le1/2\) and \(\ell\le0\), then
\(\widetilde P_{m,r}(q,\ell)\ge0\).
\end{lemma}

\begin{proof}
The identity
\begin{align}
 \rho_{n,m}(u)
 ={}&\frac mn(1-u)
       \bigl((1-u)^{n-1}-u^{n-1}\bigr)
    +\left(\frac mn-1\right)u^{n-1}
 \label{eq:dense-rho-pointwise}
\end{align}
shows that \(\rho_{n,m}(u)\ge0\) whenever \(u\le1/2\): the exponent
\(n-1\) is even and \(|1-u|\ge|u|\).  If \(\ell\le0\), then for
\(x>0\),
\[
       q+\ell\frac{1-x}{x}\le q\le\frac12.
\]
The integrand in \eqref{eq:dense-beta-second} is therefore nonnegative.
\end{proof}

Only \(\ell>0\) remains.  Pointwise positivity of \(\rho_{n,m}(q+s)\) fails
for large \(s\), so the beta integral must be used as a genuinely averaged
object.

\subsection{Laplace--gamma smoothing}

Let \(Y\sim\GammaLaw(r,1)\), with density
\[
       \frac1{\Gamma(r)}y^{r-1}e^{-y}\mathbf1_{(0,\infty)}(y).
\]

\begin{lemma}[Gamma smoothing]\label{lem:dense-gamma-smoothing}
If \(\ell>0\), then
\begin{equation}\label{eq:dense-gamma-smoothing}
 \widetilde P_{m,r}(q,\ell)
 =\frac{\ell^{n+r}}{r\Gamma(n)\Gamma(r)}
  \int_0^\infty
       t^{n+r-1}e^{-\ell t}
       \E\rho_{n,m}\left(q+\frac Yt\right)\dd t.
\end{equation}
\end{lemma}

\begin{proof}
Insert the Laplace representation
\[
       \frac1{(\ell+s)^m}
       =\frac1{\Gamma(m)}\int_0^\infty
          t^{m-1}e^{-t(\ell+s)}\dd t
\]
into \eqref{eq:dense-beta-integral}.  The exchange of integrals is justified
by absolute convergence: because \(n\) is odd, \(\rho_{n,m}\) has degree at
most \(n-1\), while \(m=n+2r\).  With \(y=ts\),
\[
 \int_0^\infty s^{r-1}e^{-ts}\rho_{n,m}(q+s)\dd s
 =\Gamma(r)t^{-r}
  \E\rho_{n,m}\left(q+\frac Yt\right).
\]
Substitution gives \eqref{eq:dense-gamma-smoothing}.
\end{proof}

The dependence on the diagonal parameter \(\ell\) now occurs only through a
positive weight.  It remains to prove a statement independent of \(\ell\).

\subsection{The shifted-gamma moment inequality}

\begin{lemma}[Uniform odd-to-even gamma moment bound]\label{lem:gamma-moment}
Let \(r,j\ge1\) be integers and \(Y\sim\GammaLaw(r,1)\).  For every
\(z\ge0\),
\begin{equation}\label{eq:gamma-moment}
       3j\,\E(zY-1)^{2j-1}
       \le(r+j)\,\E(zY-1)^{2j}.
\end{equation}
\end{lemma}

We postpone its proof for one subsection and first show exactly how it
finishes the dense region.

\begin{proposition}[Shifted-gamma positivity]\label{prop:dense-gamma-positive}
Let \(n\ge1\) be odd, \(r\ge1\), and \(m=n+2r\).  If \(0\le q\le1/3\), then
for every \(x\ge0\),
\begin{equation}\label{eq:dense-shifted-positive}
       \E\rho_{n,m}(q+xY)\ge0,
       \qquad Y\sim\GammaLaw(r,1).
\end{equation}
\end{proposition}

\begin{proof}
Set
\[
       A_n(u)=u^n+(1-u)^n.
\]
Since \(m=n+2r\), direct differentiation gives
\begin{equation}\label{eq:dense-rho-derivative}
       n\rho_{n,m}(u)=2rA_n(u)-(1-u)A_n'(u).
\end{equation}
Because \(n\) is odd,
\begin{equation}\label{eq:dense-even-expansion}
 A_n\left(\frac12+v\right)
 =\sum_{j=0}^{(n-1)/2}a_jv^{2j},
 \qquad
 a_j=2^{2j+1-n}\binom n{2j}>0.
\end{equation}
Substituting into \eqref{eq:dense-rho-derivative} yields
\begin{equation}\label{eq:dense-rho-centered}
 n\rho_{n,m}\left(\frac12+v\right)
 =2ra_0+
  \sum_{j=1}^{(n-1)/2}
  a_j\bigl(2(r+j)v^{2j}-jv^{2j-1}\bigr).
\end{equation}

Put
\[
       \delta=\frac12-q\ge\frac16.
\]
For fixed \(x\ge0\), write \(z=x/\delta\).  Then
\[
       q+xY-\frac12=\delta(zY-1).
\]
By \Cref{lem:gamma-moment}, for each \(j\ge1\),
\begin{align*}
 &\E\left[
  2(r+j)\bigl(\delta(zY-1)\bigr)^{2j}
  -j\bigl(\delta(zY-1)\bigr)^{2j-1}
 \right]\\
 &\quad\ge
 \delta^{2j-1}(r+j)
 \left(2\delta-\frac13\right)
 \E(zY-1)^{2j}\ge0.
\end{align*}
Every coefficient \(a_j\) is positive, and the constant term in
\eqref{eq:dense-rho-centered} is \(2ra_0>0\).  Taking expectations proves
\eqref{eq:dense-shifted-positive}.
\end{proof}

\begin{remark}[The threshold is retained exactly]\label{rem:dense-threshold}
The density assumption enters only through
\[
       2\delta-\frac13=2\left(\frac13-q\right).
\]
Thus gamma smoothing reaches precisely \(q=1/3\); it does not merely yield a
weaker artificial threshold.
\end{remark}

\subsection{Proof of the gamma moment inequality}\label{subsec:gamma-moment-proof}

\begin{proof}[Proof of \Cref{lem:gamma-moment}]
The case \(z=0\) is immediate.  Assume \(z>0\), put \(b=1/z\), and define
\[
       F(b)=\E(Y-b)^{2j}.
\]
Since \(F'(b)=-2j\E(Y-b)^{2j-1}\), inequality
\eqref{eq:gamma-moment} is equivalent to
\begin{equation}\label{eq:gamma-H}
       H(b):=(r+j)F(b)+\frac{3b}{2}F'(b)\ge0.
\end{equation}

For \(I_k(b)=\E(Y-b)^k\), gamma integration by parts gives
\begin{equation}\label{eq:gamma-recurrence}
       I_{k+1}(b)=(r+k-b)I_k(b)+kbI_{k-1}(b).
\end{equation}
Indeed, applying
\begin{equation}\label{eq:gamma-Stein}
       \E\bigl[Yf'(Y)\bigr]=\E\bigl[(Y-r)f(Y)\bigr]
\end{equation}
to \(f(y)=(y-b)^k\) yields \eqref{eq:gamma-recurrence}.  Taking
\(k=2j-1\), and using \(F'=-2jI_{2j-1}\) and
\(F''=2j(2j-1)I_{2j-2}\), gives the differential equation
\begin{equation}\label{eq:gamma-ODE}
       bF''(b)+(b-r-2j+1)F'(b)-2jF(b)=0.
\end{equation}

At a zero of \(H\), equations \eqref{eq:gamma-H} and
\eqref{eq:gamma-ODE} give
\begin{equation}\label{eq:gamma-crossing}
 \frac{H'(b)}{F(b)}
 =\frac{3(r+4j)b-(r+j)(5r+8j)}{3b}.
\end{equation}
Set
\begin{equation}\label{eq:gamma-bstar}
       b_*:=\frac{(r+j)(5r+8j)}{3(r+4j)}.
\end{equation}
Now \(H(0)=(r+j)\E Y^{2j}>0\), while
\[
       H(b)=(r+4j)b^{2j}+O(b^{2j-1})
       \qquad(b\to\infty).
\]
Formula \eqref{eq:gamma-crossing} says that every zero below \(b_*\) is a
strict downward crossing and every zero above \(b_*\) is a strict upward
crossing.  Therefore, if \(H\) were ever negative, its negative interval
would contain \(b_*\).  It is enough to prove
\begin{equation}\label{eq:gamma-Hbstar}
       H(b_*)>0.
\end{equation}

For the rest of the proof put \(b=b_*\), and define
\begin{equation}\label{eq:gamma-cd}
       d=\frac{2(r+j)}3,
       \qquad
       c=\frac{3jb}{r+j}.
\end{equation}
Direct calculation gives
\begin{equation}\label{eq:gamma-cd-identities}
 b-d=\frac{r(r+j)}{r+4j}>0,
 \qquad cd=2jb,
 \qquad c-d=r+2j-b.
\end{equation}
Let
\begin{equation}\label{eq:gamma-G}
       G(x)=x^{2j}(x+b)^re^{-x},
       \qquad -b\le x<\infty.
\end{equation}
Using \eqref{eq:gamma-cd-identities}, its derivative factors as
\begin{equation}\label{eq:gamma-Gprime}
 G'(x)=x^{2j-1}(x+b)^{r-1}e^{-x}(c-x)(x+d).
\end{equation}
After the change of variables \(x=y-b\), equation
\eqref{eq:gamma-H} becomes
\begin{align}
 \frac{\Gamma(r)e^b}{r+j}H(b)
 &=\int_{-b}^{\infty}
   x^{2j-1}(x-c)(x+b)^{r-1}e^{-x}\dd x\notag\\
 &=-\int_{-b}^{\infty}\frac{G'(x)}{x+d}\dd x.
 \label{eq:gamma-H-integral}
\end{align}
The apparent singularity at \(x=-d\) is removable because
\(G'(-d)=0\).  Write \(G_0=G(-d)\).  Integrating by parts separately on
the two sides of \(-d\) gives
\begin{equation}\label{eq:gamma-ibp-max}
 -\int_{-b}^{\infty}\frac{G'(x)}{x+d}\dd x
 =\frac{G_0}{b-d}
  -\int_{-b}^{\infty}\frac{G(x)-G_0}{(x+d)^2}\dd x.
\end{equation}
Thus \eqref{eq:gamma-Hbstar} follows if \(G(-d)\) is the global maximum of
\(G\).

By \eqref{eq:gamma-Gprime}, the two nonzero local maxima occur at \(-d\) and
\(c\).  Their ratio is
\begin{equation}\label{eq:gamma-max-ratio}
 \frac{G(c)}{G(-d)}
 =\left(\frac cd\right)^{2j}
  \left(\frac{b+c}{b-d}\right)^r e^{-(c+d)}.
\end{equation}
Put \(t=j/r>0\).  The logarithm of \eqref{eq:gamma-max-ratio}, divided by
\(r\), is
\begin{align}
 L(t)={}&2t\log\left(
 \frac{3t(5+8t)}{2(1+t)(1+4t)}\right)
 +\log\left(
 \frac{(1+4t)(5+8t)}{3(1+t)}\right)\notag\\
 &-\frac{32t^2+25t+2}{3(1+4t)}.
 \label{eq:gamma-L}
\end{align}
We prove
\begin{equation}\label{eq:gamma-Lnegative}
       L(t)<0\qquad(t>0).
\end{equation}
Let
\[
       A(t)=\frac{3t(5+8t)}{2(1+t)(1+4t)}.
\]
Direct differentiation gives
\begin{align}
 L'(t)
 &=2\log A(t)
 -\frac{(32t^2+16t-7)(32t^2+25t+2)}
 {3(t+1)(4t+1)^2(8t+5)},\label{eq:gamma-Lprime}\\
 L''(t)
 &=-\frac{(32t^2+25t+2)
 (128t^4+224t^3+120t^2-28t-25)}
 {t(t+1)^2(4t+1)^3(8t+5)^2}.
 \label{eq:gamma-Lsecond}
\end{align}
The quartic in \eqref{eq:gamma-Lsecond} has exactly one positive root: its
coefficient sequence has one sign change, its value at zero is negative, and
its leading coefficient is positive.  Hence \(L'\) first increases and then
decreases.  Also \(L'(t)\to-\infty\) as \(t\downarrow0\).  At \(t=1\), the
bound
\[
       \log x\ge\frac{2(x-1)}{x+1}\qquad(x\ge1)
\]
gives
\[
       L'(1)\ge\frac{76}{59}-\frac{2419}{1950}>0.
\]
At \(t=3/2\), concavity of \(\log\), together with \(\log2<7/10\), gives
\[
       L'\left(\frac32\right)
       <\frac{111}{70}-\frac{19847}{12495}<0.
\]
Thus the only interior local maximum of \(L\) lies in \((1,3/2)\), while
\[
       \lim_{t\downarrow0}L(t)=\log\frac53-\frac23<0.
\]
It remains to control \(L\) on \([1,3/2]\).  Put
\[
       B(t)=\frac{(1+4t)(5+8t)}{3(1+t)}.
\]
By concavity of \(\log\),
\[
 \log A\le\log2+\frac{A-2}{2},
 \qquad
 \log B\le\log13+\frac{B-13}{13}.
\]
Using \(\log2<7/10\) and \(\log13<13/5\), exact simplification gives
\begin{equation}\label{eq:gamma-L-rational}
 L(t)<
 \frac{3(96t^3-191t^2-16t+46)}
      {130(t+1)(4t+1)}.
\end{equation}
The cubic numerator is negative on \([1,3/2]\).  Its derivative has exactly
one zero in that interval; this critical point is a minimum, while the
endpoint values are
\[
       -65,
       \qquad -\frac{335}{4}.
\]
This proves \eqref{eq:gamma-Lnegative}.  The elementary logarithm bounds
follow directly from the exponential series.

It follows from \eqref{eq:gamma-max-ratio} that \(G(c)<G(-d)\), so
\(G(-d)\) is the global maximum.  Equation \eqref{eq:gamma-ibp-max} gives
\(H(b_*)>0\).  The zero-crossing argument then proves \(H(b)\ge0\) for all
\(b>0\), which is equivalent to \eqref{eq:gamma-moment}.
\end{proof}

\subsection{Completion of the dense-region proof}

\begin{theorem}[Dense-region theorem]\label{thm:dense-region}
If \(p\ge2/3\), then \eqref{eq:main-goal} holds for every odd \(m\ge3\).
\end{theorem}

\begin{proof}
Equivalently, \(q\le1/3\).  By
\Cref{prop:dense-expansion,lem:dense-dirichlet}, it is enough to prove
\eqref{eq:dense-diagonal-target}.  If \(\ell\le0\), this is
\Cref{lem:dense-ell-negative}.  If \(\ell>0\),
\Cref{prop:dense-gamma-positive} makes every expectation in
\eqref{eq:dense-gamma-smoothing} nonnegative, and the remaining weight is
positive.  Hence \(\widetilde P_{m,r}(q,\ell)\ge0\) for every admissible
\(r,\ell\).  Therefore every integrand in \eqref{eq:dense-expansion} is
nonnegative, so \(\Phi_m\ge0\).  Finally,
\eqref{eq:dense-defect}, \eqref{eq:dense-cycle-path}, and \eqref{eq:dense-Phi} give
\[
\begin{aligned}
 t(C_m,W)-\bigl(p^m-pq^{m-1}\bigr)
 &=q^{m-1}+S_m-t(C_m,U)\\
 &\ge q^{m-1}+S_m-x_{m-1}=\Phi_m\ge0.
\end{aligned}
\]
The finite-rank reduction returns the conclusion to arbitrary graphons.
\end{proof}

\section{The intermediate region: the one-frontier reduction for \texorpdfstring{\(m\ge9\)}{m>=9}}
\label{sec:regionII-operator}

From now on,
\begin{equation}\label{eq:regionII}
       \frac13<q<\frac12,
       \qquad \frac12<p<\frac23,
\end{equation}
and \(m\ge9\) is odd.  We continue with the common block decomposition
\eqref{eq:block-U}, the square budget \eqref{eq:HS-budget}, and the spectral
bound \eqref{eq:A-half}.  By \Cref{lem:step-reduction}, we may again work
with a step graphon during determinant manipulations.
\subsection{The one-sided spectral-shift identity}\label{subsec:shift}

The graphon operator of \(W=1-U\) has block form
\[
       T_W=
       \begin{pmatrix}
          p&-g^*\\
          -g&-A
       \end{pmatrix}.
\]
Conjugating by the sign flip on \(\one^\perp\) does not change traces, so we
may instead use
\[
       M=
       \begin{pmatrix}
          p&g^*\\
          g&-A
       \end{pmatrix}.
\]

\begin{lemma}[One-sided shift]\label{lem:shift}
For every odd \(m\ge3\),
\begin{equation}\label{eq:shift}
 t(C_m,W)=p^m-\Tr(A^m)+m[z^m]\cL(z),
\end{equation}
where
\begin{equation}\label{eq:L-series}
 \cL(z)=-\log\left(
 1-\frac{z^2\ip g{(I+zA)^{-1}g}}{1-pz}
 \right).
\end{equation}
Moreover, every coefficient of \(\cL(z)\) is nonnegative.
\end{lemma}

\begin{proof}
Schur complementation gives
\begin{align*}
 \det(I-zM)
 &=\det(I+zA)
   \left(1-pz-z^2\ip g{(I+zA)^{-1}g}\right)\\
 &=\det(I+zA)(1-pz)
   \left(1-\frac{z^2\ip g{(I+zA)^{-1}g}}{1-pz}\right).
\end{align*}
Taking \(-\log\), and using
\[
 -\log\det(I-zM)=\sum_{j\ge1}\frac{z^j}{j}\Tr(M^j),
\]
we obtain \eqref{eq:shift} by comparing the coefficient of odd degree
\(m\).

It remains to check coefficientwise nonnegativity.  Let \(\mu_g\) be the
spectral measure of \(A\) associated with \(g\).  For every integer
\(r\ge2\),
\begin{align}
 &[z^r]\frac{z^2\ip g{(I+zA)^{-1}g}}{1-pz}\notag\\
 &\qquad=\int
 \frac{p^{r-1}-(-\lambda)^{r-1}}{p+\lambda}\dd\mu_g(\lambda).
 \label{eq:X-coeff}
\end{align}
By \Cref{lem:half-norm}, \(|\lambda|\le1/2<p\).  Thus the denominator and
numerator in \eqref{eq:X-coeff} are positive.  The inner series in
\eqref{eq:L-series} therefore has nonnegative coefficients, and so does
\(-\log(1-X)=\sum_{k\ge1}X^k/k\).
\end{proof}

For odd \(m\), the linear term of the logarithm supplies the useful lower
bound
\begin{equation}\label{eq:linear-shift-lower}
 m[z^m]\cL(z)
 \ge m\int k_m(\lambda)\dd\mu_g(\lambda),
 \qquad
 k_m(\lambda)=\frac{p^{m-1}-\lambda^{m-1}}{p+\lambda}.
\end{equation}
The function \(k_m\) is strictly decreasing on \([0,p)\).  Indeed, with
\(n=m-1\),
\[
 k_m'(\lambda)
 =-\frac{p^n+np\lambda^{n-1}+(n-1)\lambda^n}{(p+\lambda)^2}<0.
\]
Also, for \(0\le s<p\),
\[
       k_m(-s)=\frac{p^{m-1}-s^{m-1}}{p-s}
       >\frac{p^{m-1}-s^{m-1}}{p+s}=k_m(s).
\]
Consequently,
\begin{equation}\label{eq:k-safe}
       k_m(\lambda)\ge k_m(L)
       \qquad\text{whenever }|\lambda|\le L<p.
\end{equation}

\subsection{No frontier, and uniqueness of a frontier}

\begin{lemma}[No-frontier case]\label{lem:no-frontier}
If the largest eigenvalue of \(A\) is at most \(q\), then
\eqref{eq:main-goal} holds.
\end{lemma}

\begin{proof}
For every eigenvalue \(\lambda\) of \(A\),
\[
       \lambda^m\le q^{m-2}\lambda^2:
\]
this is clear for \(0\le\lambda\le q\), while the left side is nonpositive
for \(\lambda<0\).  Therefore, by \Cref{lem:HS-budget},
\[
       \Tr(A^m)\le q^{m-2}\Tr(A^2)\le pq^{m-1}.
\]
The last term in \eqref{eq:shift} is nonnegative, so
\[
       t(C_m,W)\ge p^m-pq^{m-1}.
\]
\end{proof}

We henceforth assume that \(A\) has an eigenvalue
\begin{equation}\label{eq:alpha-frontier}
       \alpha>q.
\end{equation}
It is unique.  Indeed, two such eigenvalues would contribute more than
\(2q^2\) to \(\Tr(A^2)\), but
\[
       2q^2>q(1-q)=pq
\]
precisely because \(q>1/3\).  Define
\begin{equation}\label{eq:L-def}
       L=\sqrt{pq-\alpha^2}.
\end{equation}
The square budget implies that every other eigenvalue has absolute value at
most \(L\).  Moreover,
\begin{equation}\label{eq:L-order}
       0\le L<q<\alpha<\frac12<p.
\end{equation}
The strict inequality \(L>0\) will follow from the forced-variance lemma
below; no argument before that lemma requires it.
For example,
\[
 L^2<pq-q^2=q(1-2q)<q^2
\]
because \(q>1/3\), while \(\alpha<1/2\) follows from
\Cref{lem:half-norm}.

\subsection{The forced-variance ceiling}\label{subsec:forced-variance}

The next inequality restricts the admissible pair \((q,\alpha)\).

\begin{lemma}[Forced variance]\label{lem:forced-variance}
Let \(\phi\perp\one\) be a unit eigenvector with \(A\phi=\alpha\phi\).  Then
\begin{equation}\label{eq:forced-variance}
       \norm g^2\ge
       \frac{\alpha(\alpha-q)^2}{2(1-2\alpha)}.
\end{equation}
Consequently,
\begin{equation}\label{eq:alpha-ceiling-poly}
       \alpha^2+q\alpha-q\le0,
\end{equation}
and hence
\begin{equation}\label{eq:alpha-ceiling}
       q<\alpha\le r(q):=
       \frac{\sqrt{q^2+4q}-q}{2}.
\end{equation}
\end{lemma}

\begin{proof}
Put \(a_\phi=\int|\phi|\).  Since \(\int\phi=0\),
\(\int\phi_+=\int\phi_-=a_\phi/2\).  As \(0\le U\le1\),
\[
\begin{aligned}
 \alpha=\ip\phi{T_U\phi}
 &\le\iint U(x,y)
       \bigl(\phi_+(x)\phi_+(y)+\phi_-(x)\phi_-(y)\bigr)\dd x\dd y\\
 &\le\frac{a_\phi^2}{2}.
\end{aligned}
\]
Thus \(a_\phi^2\ge2\alpha\).

Write \(|\phi|=a_\phi\one+h\), where \(h\perp\one\) and
\(\norm h^2=1-a_\phi^2\).  Since \(U\ge0\),
\[
       \ip{|\phi|}{T_U|\phi|}\ge\ip\phi{T_U\phi}=\alpha.
\]
Since \(\alpha\) is the largest eigenvalue of \(A\),
\[
\begin{aligned}
 \alpha
 &\le qa_\phi^2+2a_\phi\ip g h+\ip h{Ah}\\
 &\le qa_\phi^2+2a_\phi\norm g\sqrt{1-a_\phi^2}
       +\alpha(1-a_\phi^2).
\end{aligned}
\]
Therefore
\[
       (\alpha-q)a_\phi
       \le2\norm g\sqrt{1-a_\phi^2}.
\]
The function \(t/(1-t)\) is increasing, and
\(a_\phi^2\ge2\alpha\); after squaring we obtain
\eqref{eq:forced-variance}.

Combining \eqref{eq:forced-variance} with
\(\alpha^2+2\norm g^2\le pq\), and multiplying by
\(1-2\alpha>0\), gives
\[
       (1-\alpha-q)(\alpha^2+q\alpha-q)\le0.
\]
The first factor is positive by \eqref{eq:L-order}, so
\eqref{eq:alpha-ceiling-poly} follows.  Solving the quadratic gives
\eqref{eq:alpha-ceiling}.  Finally, because \(\alpha>q\), the right side of
\eqref{eq:forced-variance} is strictly positive.  Thus \(g\ne0\), and the
square budget gives \(\alpha^2<pq\).  Consequently \(L>0\).
\end{proof}

\subsection{The master defect inequality}\label{subsec:master-defect}

Choose the sign of \(\phi\) later as convenient, and decompose
\begin{equation}\label{eq:g-decomp}
       g=\gamma\phi+g_s,
       \qquad g_s\perp\phi.
\end{equation}
For odd \(m\), define
\begin{align}
 A_m&=2L^{m-2}+mk_m(\alpha),\label{eq:Am-def}\\
 B_m&=2L^{m-2}+mk_m(L),\label{eq:Bm-def}\\
 R_m&=\alpha^m+L^m-pq^{m-1}.
 \label{eq:Rm-def}
\end{align}
By \eqref{eq:L-order}, \eqref{eq:k-safe}, and the strict decrease of
\(k_m\),
\begin{equation}\label{eq:AB-order}
       0<A_m<B_m.
\end{equation}

\begin{proposition}[Master defect]\label{prop:master-defect}
One has
\begin{equation}\label{eq:master-defect}
\begin{aligned}
 t(C_m,W)-\bigl(p^m-pq^{m-1}\bigr)
 \ge -R_m+A_m\gamma^2+B_m\norm{g_s}^2.
\end{aligned}
\end{equation}
\end{proposition}

\begin{proof}
Let \((\lambda_i)\) be the safe eigenvalues, i.e. all eigenvalues of \(A\)
other than \(\alpha\).  Since \(|\lambda_i|\le L\) and \(m\) is odd,
\[
       \sum_i\lambda_i^m
       \le L^{m-2}\sum_i\lambda_i^2.
\]
The budget \eqref{eq:HS-budget} and \eqref{eq:g-decomp} give
\[
       \sum_i\lambda_i^2
       \le L^2-2\gamma^2-2\norm{g_s}^2.
\]
Hence
\begin{equation}\label{eq:trace-payment}
 -\Tr(A^m)
 \ge-\alpha^m-L^m
    +2L^{m-2}\bigl(\gamma^2+\norm{g_s}^2\bigr).
\end{equation}
By \eqref{eq:linear-shift-lower} and \eqref{eq:k-safe},
\begin{equation}\label{eq:shift-payment}
 m[z^m]\cL(z)
 \ge mk_m(\alpha)\gamma^2+mk_m(L)\norm{g_s}^2.
\end{equation}
Combining \eqref{eq:shift}, \eqref{eq:trace-payment}, and
\eqref{eq:shift-payment} yields \eqref{eq:master-defect}.
\end{proof}

\section{Eliminating the frontier eigenfunction: a Huber envelope}
\label{sec:huber}

The terminology refers to the same quadratic-versus-linear active-set
geometry as Huber's robust loss~\cite{Huber1964}.  The exact dual formula
below is a compact convex--concave minimax interchange of the type developed
by Sion~\cite{Sion1958}; we nevertheless prove the required special case
directly.

The master defect still contains the coupling variables \(\gamma\) and
\(g_s\).  We now eliminate them together with the shape of the frontier
eigenfunction.

Choose the sign of \(\phi\) so that
\[
       b:=\ip{|\phi|}{\phi}\ge0.
\]
Set
\begin{equation}\label{eq:shape-defs}
 z=\left(\int|\phi|\right)^2,
 \qquad
 k=|\phi|-\sqrt z\,\one-b\phi,
 \qquad
 K=\norm k^2.
\end{equation}
Then
\begin{equation}\label{eq:shape-budget}
       z+b^2+K=1,
       \qquad z\ge2\alpha.
\end{equation}

\subsection{The direct and safe channels}

\begin{lemma}[Direct channel]\label{lem:direct-channel}
The frontier coupling obeys
\begin{equation}\label{eq:direct-channel}
       \gamma\le u_\gamma:=
       \frac{z-2\alpha-2\alpha b}{2\sqrt z}.
\end{equation}
\end{lemma}

\begin{proof}
Since \(T_U\phi=\gamma\one+\alpha\phi\) and
\[
       T_U\phi\le T_U\phi_+\le\int\phi_+=\frac{\sqrt z}{2}
\]
pointwise, pairing with \(\phi_+\) gives
\[
       \frac{\gamma\sqrt z}{2}
       +\frac{\alpha(1+b)}{2}
       \le\frac z4.
\]
This is \eqref{eq:direct-channel}.
\end{proof}

\begin{lemma}[Safe channel]\label{lem:safe-channel}
Put
\begin{equation}\label{eq:H-shape}
       \cH=\frac{(\alpha-q)z+(\alpha-L)K}{2\sqrt z}.
\end{equation}
Then
\begin{equation}\label{eq:safe-channel}
       b\gamma+\ip{g_s}{k}\ge\cH.
\end{equation}
Consequently, if \(K>0\),
\begin{equation}\label{eq:gs-safe-lower}
       \norm{g_s}^2\ge
       \frac{\pos{\cH-b\gamma}^2}{K}.
\end{equation}
\end{lemma}

\begin{proof}
Because \(U\ge0\),
\[
 \ip{|\phi|}{T_U|\phi|}-\ip\phi{T_U\phi}
 =4\ip{\phi_+}{T_U\phi_-}\ge0.
\]
Expanding \(|\phi|=\sqrt z\,\one+b\phi+k\), using
\(k\perp\one,\phi\), and using that every safe eigenvalue is at most \(L\),
we obtain
\[
\begin{aligned}
 \alpha
 &\le qz+2\sqrt z\bigl(b\gamma+\ip{g_s}{k}\bigr)
       +\alpha b^2+\ip k{Ak}\\
 &\le qz+2\sqrt z\bigl(b\gamma+\ip{g_s}{k}\bigr)
       +\alpha b^2+LK.
\end{aligned}
\]
Since \(z+b^2+K=1\), this rearranges to
\eqref{eq:safe-channel}.  If \(\cH-b\gamma>0\), Cauchy--Schwarz gives
\(\norm{g_s}\sqrt K\ge\cH-b\gamma\); otherwise
\eqref{eq:gs-safe-lower} is trivial.
\end{proof}

For fixed shape variables, \Cref{prop:master-defect} and
\Cref{lem:direct-channel,lem:safe-channel} reduce the payment to
\begin{equation}\label{eq:Q-shape}
 \cQ=
 \min_{\gamma\le u_\gamma}
 \left\{
 A_m\gamma^2+\frac{B_m}{K}\pos{\cH-b\gamma}^2
 \right\},
\end{equation}
with the natural limiting interpretation when \(K=0\).

\subsection{Exact shape elimination}\label{subsec:shape-elim}

Define
\begin{equation}\label{eq:dfe}
       d=\alpha-q,
       \qquad f=\alpha-L,
       \qquad e=1-2\alpha,
\end{equation}
so \(d,f,e>0\).  Put
\begin{equation}\label{eq:C-xi-rho}
 C_m=\frac{B_mf\sqrt{2\alpha}\,e^2}{4\alpha^2},
 \qquad
 \xi=\frac{4\alpha^2d}{e^2},
 \qquad
 \rho=\frac{A_m}{B_m}\frac{\sqrt\alpha}{2\sqrt2\,f},
\end{equation}
and define
\begin{equation}\label{eq:psi-def}
       \psi(\xi,\rho)=
       \min_{0\le v\le1}
       \left\{\rho v^2+\pos{\xi-v+v^2}\right\}.
\end{equation}

\begin{theorem}[Huber shape elimination]\label{thm:huber-elim}
For every admissible frontier configuration,
\begin{equation}\label{eq:huber-payment}
       A_m\gamma^2+B_m\norm{g_s}^2
       \ge C_m\psi(\xi,\rho).
\end{equation}
Consequently, \eqref{eq:main-goal} follows once
\begin{equation}\label{eq:scalar-target}
       R_m\le C_m\psi(\xi,\rho)
\end{equation}
is proved for every pair \((q,\alpha)\) satisfying
\eqref{eq:regionII} and \eqref{eq:alpha-ceiling}.
\end{theorem}

\begin{proof}
We first treat \(K>0\).

Suppose \(u_\gamma\le0\).  Every admissible \(\gamma\) is nonpositive, so
\(\cH-b\gamma\ge\cH\).  Hence
\[
       \cQ\ge B_m\frac{\cH^2}{K}.
\]
By \eqref{eq:H-shape},
\[
 \frac{\cH^2}{K}
 =\frac{(dz+fK)^2}{4zK}\ge df,
\]
because this inequality is equivalent to \((dz-fK)^2\ge0\).  On the other
hand, taking \(v=0\) in \eqref{eq:psi-def} gives
\[
 C_m\psi(\xi,\rho)
 \le C_m\xi=B_mfd\sqrt{2\alpha}\le B_mfd.
\]
Thus \eqref{eq:huber-payment} holds in this case.

Now suppose \(u_\gamma>0\).  A minimizer in \eqref{eq:Q-shape} may be taken
with \(\gamma\ge0\): replacing a negative \(\gamma\) by \(0\) decreases both
terms.  The direct constraint gives
\[
       2\gamma\sqrt z+2\alpha b
       \le z-2\alpha=e-b^2-K\le e-b^2.
\]
Since \(z\ge2\alpha\),
\begin{equation}\label{eq:b-gamma-budget}
       b^2+2\alpha b+2\gamma\sqrt{2\alpha}\le e.
\end{equation}
It follows that
\begin{equation}\label{eq:beta-gamma}
 0\le\gamma\le\frac{e}{2\sqrt{2\alpha}},
 \qquad
 b\le\beta(\gamma):=
 \sqrt{\alpha^2+e-2\gamma\sqrt{2\alpha}}-\alpha.
\end{equation}
Rationalizing the square root,
\begin{equation}\label{eq:beta-rational}
       \beta(\gamma)
       \le\frac{e-2\gamma\sqrt{2\alpha}}{2\alpha}.
\end{equation}
Also, since \(z\ge2\alpha\) and \(z\le1\),
\begin{equation}\label{eq:H-lower}
       \cH\ge\frac{d\sqrt{2\alpha}}2+\frac{fK}{2}.
\end{equation}
For every real \(w\),
\begin{equation}\label{eq:K-AMGM}
       \inf_{K>0}\frac{\pos{w+fK/2}^2}{K}=2f\pos w.
\end{equation}
Indeed, for \(w\ge0\) this is the arithmetic--geometric mean inequality,
with equality at \(K=2w/f\); for \(w<0\), both sides vanish after choosing
\(0<K\le-2w/f\).  Using \eqref{eq:gs-safe-lower},
\eqref{eq:H-lower}, and \eqref{eq:K-AMGM}, and then maximizing \(b\) by
\eqref{eq:beta-gamma}, we obtain
\begin{equation}\label{eq:one-d-gamma}
 \cQ\ge
 \min_{0\le\gamma\le e/(2\sqrt{2\alpha})}
 \left\{
 A_m\gamma^2+2B_mf
 \pos{\frac{d\sqrt{2\alpha}}2-\gamma\beta(\gamma)}
 \right\}.
\end{equation}
Set
\[
       v=\frac{2\gamma\sqrt{2\alpha}}e\in[0,1].
\]
By \eqref{eq:beta-rational},
\[
       \gamma\beta(\gamma)
       \le\frac{e^2v(1-v)}{4\alpha\sqrt{2\alpha}}.
\]
Therefore
\[
 \frac{d\sqrt{2\alpha}}2-\gamma\beta(\gamma)
 \ge\frac{\sqrt{2\alpha}\,e^2}{8\alpha^2}
       (\xi-v+v^2).
\]
Furthermore,
\[
       A_m\gamma^2=C_m\rho v^2,
       \qquad
       2B_mf\frac{\sqrt{2\alpha}\,e^2}{8\alpha^2}=C_m.
\]
Substitution in \eqref{eq:one-d-gamma} proves
\eqref{eq:huber-payment}.

Finally suppose \(K=0\).  The safe channel becomes
\[
       b\gamma\ge\frac{d\sqrt z}{2}
       \ge\frac{d\sqrt{2\alpha}}2>0,
\]
so \(b,\gamma>0\) and \eqref{eq:b-gamma-budget} still applies.  With the
same variable \(v\),
\[
 \frac{d\sqrt{2\alpha}}2
 \le b\gamma
 \le\frac{e^2v(1-v)}{4\alpha\sqrt{2\alpha}},
\]
which is equivalent to \(\xi\le v-v^2\).  Thus the positive-part term in
\eqref{eq:psi-def} vanishes at this \(v\), and
\[
       A_m\gamma^2=C_m\rho v^2\ge C_m\psi(\xi,\rho).
\]
This completes the proof.
\end{proof}

\subsection{The dual envelope and its two canonical witnesses}

\begin{proposition}[Dual form of the Huber envelope]\label{prop:huber-dual}
For \(\xi\ge0\) and \(\rho>0\),
\begin{equation}\label{eq:psi-dual}
       \psi(\xi,\rho)=
       \max_{0\le\lambda\le1}
       \left\{
       \lambda\xi-\frac{\lambda^2}{4(\rho+\lambda)}
       \right\}.
\end{equation}
\end{proposition}

\begin{proof}
Use
\[
       \pos x=\max_{0\le\lambda\le1}\lambda x.
\]
The resulting function
\[
       (\rho+\lambda)v^2-\lambda v+\lambda\xi
\]
is convex in \(v\), affine in \(\lambda\), and the two domains are compact
convex intervals.  The elementary finite-dimensional minimax theorem (or
Sion's theorem) permits the interchange of \(\min_v\) and \(\max_\lambda\).
For fixed \(\lambda\), the minimizer is
\[
       v=\frac{\lambda}{2(\rho+\lambda)}\in[0,1/2],
\]
and substitution gives \eqref{eq:psi-dual}.
\end{proof}

Two choices of \(\lambda\) will be used throughout.

\begin{corollary}[Quadratic and linear witnesses]\label{cor:two-witnesses}
If \(2\rho\xi\le1\), then
\begin{equation}\label{eq:witness-quadratic}
 C_m\psi(\xi,\rho)
 \ge
 \frac{2\alpha^3A_md^2}{e^2}
 \frac{1+4\xi}{1+2\xi}.
\end{equation}
If \(2\rho\xi>1\), then
\begin{equation}\label{eq:witness-linear}
 C_m\psi(\xi,\rho)
 >B_mf\sqrt{2\alpha}\,d
 \frac{1+4\xi}{2+4\xi}.
\end{equation}
\end{corollary}

\begin{proof}
In the first case, take \(\lambda=2\rho\xi\) in
\eqref{eq:psi-dual}.  This gives
\[
       \psi(\xi,\rho)\ge
       \rho\xi^2\frac{1+4\xi}{1+2\xi}.
\]
A direct calculation from \eqref{eq:C-xi-rho} gives
\[
       C_m\rho\xi^2=\frac{2\alpha^3A_md^2}{e^2}.
\]

In the second case, take \(\lambda=1\).  Since
\(\rho>1/(2\xi)\),
\[
       \frac1{4(1+\rho)}<\frac{\xi}{2+4\xi}.
\]
Therefore
\[
 \psi(\xi,\rho)
 \ge\xi-\frac1{4(1+\rho)}
 >\xi\frac{1+4\xi}{2+4\xi}.
\]
Finally,
\[
       C_m\xi=B_mf\sqrt{2\alpha}\,d.
\]
\end{proof}

\section{The dimensionless chart}\label{sec:chart}

Set
\begin{equation}\label{eq:N-def}
       N=m-2.
\end{equation}
Thus \(N\ge7\) is odd.  Introduce
\begin{equation}\label{eq:chart-basic}
 a=\frac q\alpha,
 \qquad
 \ell=\frac L\alpha,
 \qquad
 u=1-a,
 \qquad
 D=1+a^2+\ell^2,
 \qquad
 \sigma=\frac p\alpha.
\end{equation}
The relation \(L^2=pq-\alpha^2\) is equivalent to
\begin{equation}\label{eq:chart-inverse-1}
       \alpha=\frac aD,
       \qquad q=\frac{a^2}{D},
       \qquad \sigma=\frac{1+\ell^2}{a}.
\end{equation}
Moreover,
\begin{equation}\label{eq:chart-d-e}
       d=\frac{au}{D},
       \qquad e=\frac{u^2+\ell^2}{D}.
\end{equation}

\begin{lemma}[Chart domain]\label{lem:chart-domain}
The admissible parameters satisfy
\begin{equation}\label{eq:chart-domain}
\begin{gathered}
 0<\ell<a<1,
 \qquad
 D<3a^2,
 \qquad
 \ell^2\ge a(1-a)=au,\\
 a>\frac{1+\sqrt{13}}6>\frac34,
 \qquad
 1<\sigma<2.
\end{gathered}
\end{equation}
\end{lemma}

\begin{proof}
The order \(0<L<q<\alpha\) gives \(0<\ell<a<1\).  Since
\(q=a^2/D>1/3\), we have \(D<3a^2\).  The ceiling
\(\alpha^2+q\alpha-q\le0\), after substitution from
\eqref{eq:chart-inverse-1}, is exactly
\[
       \ell^2\ge a-a^2=au.
\]
Hence \(D\ge1+a\).  Combining this with \(D<3a^2\) gives
\(3a^2-a-1>0\), which is the stated lower bound on \(a\).  Finally,
\(\sigma=p/\alpha>1\) because \(p>1/2>\alpha\), while
\(p<2/3\) and \(\alpha>q>1/3\) give \(\sigma<2\).
\end{proof}

Now put
\begin{equation}\label{eq:zeta-v-def}
       \zeta=\frac{\ell^2}{u},
       \qquad v=\sigma-1,
       \qquad H_\zeta=\zeta+1+v.
\end{equation}
Solving these relations gives
\begin{equation}\label{eq:zeta-v-inverse}
       u=\frac v{H_\zeta},
       \qquad
       a=\frac{\zeta+1}{H_\zeta},
       \qquad
       \ell^2=\frac{\zeta v}{H_\zeta}.
\end{equation}
The frontier ceiling \(\ell^2\ge au\) becomes
\begin{equation}\label{eq:zeta-domain}
       \zeta(\zeta+v)\ge1.
\end{equation}
Equivalently,
\begin{equation}\label{eq:zeta-simple-lower}
       \zeta\ge\frac{\sqrt{v^2+4}-v}{2}\ge1-\frac v2.
\end{equation}
In particular,
\begin{equation}\label{eq:zeta-lower}
       \zeta\ge a>\frac34,
       \qquad 0<v<1.
\end{equation}

For later use, define
\begin{equation}\label{eq:Q-chart}
       Q=v^2+v\zeta+2v+2\zeta+2.
\end{equation}
A direct substitution gives
\begin{equation}\label{eq:D-Q}
       D=\frac{(\zeta+1)Q}{H_\zeta^2}.
\end{equation}
Since \(D<3a^2\),
\begin{equation}\label{eq:Q-upper}
       Q<3(\zeta+1).
\end{equation}
We shall repeatedly use the identity
\begin{equation}\label{eq:H-identity}
 \zeta H_\zeta^2-(\zeta+1)(\zeta+v)^2
 =\zeta^2+\zeta-v^2>0.
\end{equation}
The positivity follows from \eqref{eq:zeta-domain}:
\[
 \zeta^2+\zeta-v^2
 \ge1+\zeta(1-v)-v^2
 =(1-v)(1+v+\zeta)>0.
\]

\subsection{The normalized defect and finite quotients}

Normalize \(R_m\) by \(\alpha^m\):
\begin{equation}\label{eq:F-def}
       F_N:=\frac{R_m}{\alpha^m}
       =1-a^N-\ell^2(a^N-\ell^N).
\end{equation}
Define
\begin{equation}\label{eq:T-def}
       T_N=N-\zeta(a^N-\ell^N).
\end{equation}
Since \(1-a^N\le N(1-a)=Nu\),
\begin{equation}\label{eq:F-T}
       F_N\le uT_N.
\end{equation}
If \(T_N\le0\), then \(R_m\le0\), and the scalar target
\eqref{eq:scalar-target} is immediate.  We may therefore assume
\begin{equation}\label{eq:T-positive}
       T_N>0
\end{equation}
whenever convenient.

The normalized versions of \(A_m\) and \(B_m\) are
\begin{align}
 \frac{A_m}{\alpha^N}
 &=2\ell^N+(N+2)K_A,
 &
 K_A&=\frac{\sigma^{N+1}-1}{\sigma+1},
 \label{eq:KA-def}\\
 \frac{B_m}{\alpha^N}
 &=2\ell^N+(N+2)K_L,
 &
 K_L&=\frac{\sigma^{N+1}-\ell^{N+1}}{\sigma+\ell}.
 \label{eq:KL-def}
\end{align}
Because \(N+1\) is even,
\begin{align}
 K_A&=(\sigma-1)
       \sum_{j=0}^{(N-1)/2}\sigma^{2j},
 \label{eq:KA-sum}\\
 K_L&=(\sigma-\ell)
       \sum_{j=0}^{(N-1)/2}
       \sigma^{N-1-2j}\ell^{2j}.
 \label{eq:KL-sum}
\end{align}

\subsection{Three elementary bounds for \texorpdfstring{\(T_N\)}{TN}}

\begin{lemma}[Defect bounds]\label{lem:T-bounds}
For every admissible chart point and every odd \(N\ge7\),
\begin{align}
 T_N&\le N(1+v)-\zeta(1-v^{N/2}),
 \label{eq:T-bound-one}\\
 T_N&<N(1-\ell^{N+1}).
 \label{eq:T-bound-two}
\end{align}
If \(T_N>0\), then also
\begin{equation}\label{eq:T-tangent-bound}
       \zeta T_N<\frac{3N(N+1)}2(1+v).
\end{equation}
\end{lemma}

\begin{proof}
By Bernoulli's inequality,
\[
       a^N\ge1-Nu.
\]
Also \(\ell^2=\zeta v/H_\zeta<v\), so
\(\ell^N<v^{N/2}\).  Since \(\zeta u< v\),
\[
\begin{aligned}
 T_N
 &=N-\zeta(a^N-\ell^N)\\
 &\le N-\zeta+N\zeta u+\zeta v^{N/2}\\
 &\le N(1+v)-\zeta(1-v^{N/2}),
\end{aligned}
\]
proving \eqref{eq:T-bound-one}.

Next, \(D<3a^2\) gives
\[
       \ell^2<2a^2-1<(2a-1)^2.
\]
Thus \(a-\ell>1-a=u\).  By convexity of \(x\mapsto x^N\),
\[
       a^N-\ell^N>N(a-\ell)\ell^{N-1}>Nu\ell^{N-1}.
\]
Multiplication by \(\zeta=\ell^2/u\) gives
\eqref{eq:T-bound-two}.

Assume now \(T_N>0\).  From \eqref{eq:T-bound-one},
\begin{equation}\label{eq:zeta-upper-Z}
       \zeta<Z:=\frac{N(1+v)}{1-v^{N/2}}.
\end{equation}
Put \(r=(N+1)/2\) and
\[
       g(t)=t\left[1-
       \left(\frac{tv}{t+1+v}\right)^r\right].
\]
If \(\theta=(1+v)/(t+1+v)\), then
\[
       g'(t)=1-[v(1-\theta)]^r(1+r\theta)\ge0,
\]
because \(v^r\le1\) and Bernoulli's inequality gives
\((1-\theta)^{-r}\ge1+r\theta\).  Therefore,
using \eqref{eq:T-bound-two} and \eqref{eq:zeta-upper-Z},
\begin{equation}\label{eq:gZ-step}
 \zeta T_N<N g(\zeta)<Ng(Z).
\end{equation}
Let
\[
       \eta=v^{N/2},
       \qquad
       \omega=\frac{Nv}{N+1-\eta}.
\]
A calculation gives
\[
       Ng(Z)=
       \frac{N^2(1+v)}{1-\eta}
       \left(1-\omega^{(N+1)/2}\right).
\]
As a function of \(\eta\in(0,1)\),
\[
 f(\eta)=\omega^{(N+1)/2}
 =\left(\frac{N}{N+1-\eta}\right)^{(N+1)/2}
  \eta^{(N+1)/N}
\]
is convex.  Indeed, with
\(r=(N+1)/2\), \(s=(N+1)/N\), and \(d=N+1-\eta\),
\[
       \frac{f''(\eta)}{f(\eta)}
       =\frac{r(r+1)}{d^2}
        +\frac{2rs}{d\eta}
        +\frac{s(s-1)}{\eta^2}>0.
\]
Moreover,
\[
       f(1)=1,
       \qquad
       f'(1)=\frac{3(N+1)}{2N}.
\]
The tangent inequality for a convex function therefore yields
\[
       1-f(\eta)
       \le\frac{3(N+1)}{2N}(1-\eta).
\]
Substitution in \eqref{eq:gZ-step} proves
\eqref{eq:T-tangent-bound}.
\end{proof}

\section{The quadratic Huber branch}\label{sec:quadratic}

Assume
\begin{equation}\label{eq:quadratic-branch}
       2\rho\xi\le1.
\end{equation}
We shall prove \eqref{eq:scalar-target} for every odd \(N\ge7\), hence for
every odd \(m\ge9\).

By \eqref{eq:witness-quadratic}, discarding the factor
\((1+4\xi)/(1+2\xi)>1\), and using \eqref{eq:KA-def},
\begin{equation}\label{eq:quadratic-payment-start}
 \frac{C_m\psi(\xi,\rho)}{\alpha^m}
 \ge2(N+2)\alpha\left(\frac de\right)^2K_A.
\end{equation}
Since the summands in \eqref{eq:KA-sum} have geometric mean
\(\sigma^{(N-1)/2}\), AM--GM gives
\begin{equation}\label{eq:KA-AMGM}
       K_A\ge\frac{N+1}{2}
       v\sigma^{(N-1)/2}.
\end{equation}
Using \eqref{eq:chart-inverse-1}, \eqref{eq:chart-d-e}, and
\(v=u(1+\zeta)/a\), the right side of
\eqref{eq:quadratic-payment-start} is at least
\begin{equation}\label{eq:quad-payment-M}
 u(N+2)(N+1)
 \frac{a^2(1+\zeta)}{D(u+\zeta)^2}
 \sigma^{(N-1)/2}.
\end{equation}

\begin{lemma}[Coefficient estimate]\label{lem:quad-coeff}
One has
\begin{equation}\label{eq:M-lower}
       \frac{a^2(1+\zeta)}{D(u+\zeta)^2}
       >\frac1{3\zeta}.
\end{equation}
\end{lemma}

\begin{proof}
Using \eqref{eq:zeta-v-inverse} and \eqref{eq:D-Q}, the left side is
\[
       M=\frac{H_\zeta^2}{(\zeta+v)^2Q}.
\]
By \eqref{eq:H-identity},
\[
       \frac{H_\zeta^2}{(\zeta+v)^2}
       >\frac{\zeta+1}{\zeta},
\]
while \eqref{eq:Q-upper} gives \(Q<3(\zeta+1)\).  Multiplication proves
\eqref{eq:M-lower}.
\end{proof}

Combining \eqref{eq:quad-payment-M} and \eqref{eq:M-lower},
\begin{equation}\label{eq:quad-payment-final}
 \frac{C_m\psi(\xi,\rho)}{\alpha^m}
 >u\frac{(N+2)(N+1)}{3\zeta}
    (1+v)^{(N-1)/2}.
\end{equation}
In view of \eqref{eq:F-T}, it remains to pay for \(uT_N\).

\begin{proposition}[Quadratic branch]\label{prop:quadratic-branch}
For every odd \(N\ge7\), assumption \eqref{eq:quadratic-branch} implies
\[
       R_m\le C_m\psi(\xi,\rho).
\]
\end{proposition}

\begin{proof}
If \(T_N\le0\), the claim follows from \eqref{eq:F-T}.  Assume
\(T_N>0\).

First, \eqref{eq:T-bound-one} implies
\[
 \zeta T_N
 \le N(1+v)\zeta-(1-v^{N/2})\zeta^2
 \le\frac{N^2(1+v)^2}{4(1-v^{N/2})}.
\]
Thus \eqref{eq:quad-payment-final} is sufficient whenever
\begin{equation}\label{eq:E-N}
 E_N(v):=
 \frac{4(N+2)(N+1)}{3N^2}
 (1-v^{N/2})(1+v)^{(N-5)/2}\ge1.
\end{equation}
Second, the tangent estimate \eqref{eq:T-tangent-bound} shows that
\eqref{eq:quad-payment-final} is sufficient whenever
\begin{equation}\label{eq:H-N}
 H_N(v):=
 \frac{2(N+2)}{9N}(1+v)^{(N-3)/2}\ge1.
\end{equation}
We now cover the whole interval \(0<v<1\).

If \(N\ge9\) and \(v\le3/5\), then
\[
 E_N(v)\ge\frac43\left(1-\left(\frac35\right)^4\right)
 =\frac{2176}{1875}>1.
\]
If \(N\ge9\) and \(v\ge3/5\), then
\[
 \frac{H_{N+2}(v)}{H_N(v)}
 =(1+v)\frac{N(N+4)}{(N+2)^2}>1,
\]
so \(H_N(v)\) increases under \(N\mapsto N+2\), and
\[
       H_9(3/5)=\frac{11264}{10125}>1.
\]

It remains to take \(N=7\).  For \(v\le9/10\),
\[
 E_7(v)\ge\frac{96}{49}(1-v^3)(1+v)\ge1.
\]
Indeed, \(1+v-v^3-v^4\) is concave, so its minimum on
\([0,9/10]\) occurs at an endpoint; at \(v=9/10\) it equals
\(5149/10000>49/96\).  For \(v\ge9/10\),
\[
       H_7(v)\ge\frac27\left(\frac{19}{10}\right)^2
       =\frac{361}{350}>1.
\]
Thus either \eqref{eq:E-N} or \eqref{eq:H-N} holds in every case, proving
the proposition.
\end{proof}

\section{The linear Huber branch}\label{sec:linear}

Assume
\begin{equation}\label{eq:linear-branch}
       2\rho\xi>1.
\end{equation}
Abbreviate
\begin{equation}\label{eq:varphi-cxi}
       \varphi=\frac{1+4\xi}{2+4\xi},
       \qquad
       c_\xi=\sqrt{2\alpha}\,\varphi.
\end{equation}
By \eqref{eq:witness-linear}, \eqref{eq:KL-def}, and
\(f=\alpha(1-\ell)\), \(d=\alpha u\),
\begin{equation}\label{eq:linear-payment-start}
 \frac{C_m\psi(\xi,\rho)}{\alpha^m}
 >u(N+2)c_\xi(1-\ell)K_L.
\end{equation}
From \eqref{eq:KL-sum},
\begin{equation}\label{eq:KL-first-term}
       K_L\ge(\sigma-\ell)\sigma^{N-1}.
\end{equation}
In view of \eqref{eq:F-T}, it is therefore enough to prove
\begin{equation}\label{eq:linear-core}
       T_N\le
       (N+2)c_\xi(1-\ell)(\sigma-\ell)\sigma^{N-1}.
\end{equation}

\subsection{The compensation lemma}

The factor \(c_\xi\) may be small, but only where a short power of
\(\sigma\) compensates for it.

\begin{lemma}[Compensation]\label{lem:compensation}
On the whole admissible chart,
\begin{equation}\label{eq:compensation}
       c_\xi\sigma^{\zeta/4}\ge1-\ell^2.
\end{equation}
\end{lemma}

\begin{proof}
We first show
\begin{equation}\label{eq:sqrt-compensation}
       \sqrt{2\alpha}\ge1-\ell^2.
\end{equation}
Using \(\alpha=a/D\), the square of \eqref{eq:sqrt-compensation} is
\[
       2a\ge(1+a^2+y)(1-y)^2,
       \qquad y=\ell^2.
\]
The difference between the two sides has derivative
\[
       (1-y)(1+2a^2+3y)>0.
\]
The smallest admissible \(y\) is \(a(1-a)\), where the difference is
\[
       (1-a)(a^4+a^2+2a-1)>0
\]
because \(a>3/4\).  This proves \eqref{eq:sqrt-compensation}.

Next we show \(\varphi\sigma^{\zeta/4}\ge1\).  From
\eqref{eq:C-xi-rho}, \eqref{eq:zeta-v-inverse}, and
\eqref{eq:D-Q},
\begin{equation}\label{eq:xi-zeta-v}
       \xi=\frac{4H_\zeta^2}
       {v(\zeta+v)^2Q}.
\end{equation}
Equations \eqref{eq:Q-upper} and \eqref{eq:H-identity} imply
\begin{equation}\label{eq:xi-comp-bound}
       4\zeta v\xi>\frac{16}{3}.
\end{equation}
We claim that
\begin{equation}\label{eq:key-log-comp}
       \zeta v(1+4\xi)>2(v+2).
\end{equation}
Indeed, the left side is larger than \(\zeta v+16/3\).  If
\(\zeta\ge2\), this is immediate.  If \(\zeta<2\), then
\[
       v(2-\zeta)<\frac54<\frac43
\]
because \(v<1\) and \(\zeta>3/4\), which again proves
\eqref{eq:key-log-comp}.

The elementary inequality
\[
       \log(1+v)\ge\frac{2v}{v+2}
\]
now gives
\[
       \frac\zeta4\log\sigma>
       \frac1{1+4\xi}.
\]
Exponentiating and using \(e^t>1+t\),
\[
       \sigma^{\zeta/4}>
       1+\frac1{1+4\xi}
       =\frac{2+4\xi}{1+4\xi}=\frac1\varphi.
\]
Together with \eqref{eq:sqrt-compensation}, this proves
\eqref{eq:compensation}.
\end{proof}

\subsection{Two broad estimates for the linear payment}

We first dispose of the range \(\zeta\ge N\) when \(v\) is not large.

\begin{lemma}[Above the cycle scale]\label{lem:linear-high-zeta}
If \(\zeta\ge N\), \(0<v\le5/8\), and \(N\ge7\) is odd, then
\eqref{eq:linear-core} holds.
\end{lemma}

\begin{proof}
By \eqref{eq:T-bound-one} and \(\zeta\ge N\),
\begin{equation}\label{eq:T-high-zeta}
       T_N\le N(v+v^{N/2}).
\end{equation}
Put \(y=\sqrt v\).  Since \(\ell^2<v\),
\[
       (1-\ell)(\sigma-\ell)
       \ge(1-y)(1-y+y^2).
\]
Also \(q<\alpha\) and \(p=(1+v)\alpha\) imply
\(\alpha>1/(2+v)\), while \(\varphi>1/2\).  Hence
\begin{equation}\label{eq:cxi-crude}
       c_\xi>\frac1{\sqrt{2(2+v)}}.
\end{equation}
It is enough to show
\begin{equation}\label{eq:h-sufficient}
 h(y):=
 \frac{(1-y)(1-y+y^2)(1+y^2)^5}{y^2}
 >\sqrt{2(2+y^2)}\,(1+y^{N-2}).
\end{equation}
Indeed, the omitted factors
\((N+2)/N\) and \((1+y^2)^{N-6}\) are at least \(1\).

We claim that \(h\) is decreasing on \((0,4/5]\).  Differentiation gives
\[
 h'(y)=-\frac{(1+y^2)^4}{y^3}S(y),
\]
where
\[
       S(y)=11y^5-20y^4+19y^3-8y^2-2y+2.
\]
For \(0\le y\le1/2\),
\[
 S(y)\ge2-2y-8y^2+9y^3\ge\frac18.
\]
Here the first inequality follows from
\(y^3(11y^2-20y+10)\ge0\), and the cubic on the right is decreasing on this
interval.  For \(1/2\le y\le4/5\), write \(y=2/3+t\).  If \(t\ge0\), then \(0\le t\le2/15\), and
\[
\begin{aligned}
 S(y)&=\frac{58}{243}-\frac{14}{81}t
      +\frac{250}{27}t^2+\frac{131}{9}t^3
      +\frac{50}{3}t^4+11t^5\\
 &\ge \frac{58}{243}-\frac{14}{81}\frac{2}{15}
  =\frac{262}{1215}>0.
\end{aligned}
\]
If \(t=-s\le0\), where \(0\le s\le1/6\), then
\[
\begin{aligned}
 S(y)={}&\frac{58}{243}+\frac{14}{81}s
 +\frac{250}{27}s^2-\frac{131}{9}s^3
 +\frac{50}{3}s^4-11s^5\\
 \ge{}&\frac{58}{243}+\frac{14}{81}s
 +\frac{1465}{216}s^2>0.
\end{aligned}
\]
Thus \(h'\le0\).

Since \(v\le5/8\), we have \(y<4/5\), and therefore
\[
       h(y)\ge h(4/5)
       =\frac{2432980221}{781250000}>\frac{46}{15}.
\]
On the other hand, \(N\ge7\) gives
\[
 \sqrt{2(2+y^2)}(1+y^{N-2})
 \le\sqrt{\frac{21}{4}}
       \left(1+\left(\frac45\right)^5\right)
 <\frac{46}{15};
\]
after squaring, the final difference is
\(52766011/351562500>0\).  This proves
\eqref{eq:h-sufficient} and hence \eqref{eq:linear-core}.
\end{proof}

The next estimate handles all \(\zeta\) once \(v\) is large.

\begin{lemma}[Large \(v\)]\label{lem:linear-large-v}
If \(v\ge5/8\) and \(N\ge7\) is odd, then
\eqref{eq:linear-core} holds.
\end{lemma}

\begin{proof}
By \eqref{eq:T-bound-two},
\begin{equation}\label{eq:T-geometric-cancel}
       T_N<N(N+1)(1-\ell).
\end{equation}
Using \eqref{eq:cxi-crude}, \(\ell\le\sqrt v\), and
\(\sigma=1+v\), it is enough, after cancelling \(1-\ell\), to prove
\begin{equation}\label{eq:large-v-target}
 \frac{N+2}{\sqrt{2(2+v)}}
 (1+v-\sqrt v)(1+v)^{N-1}
 \ge N(N+1).
\end{equation}
The left side is increasing in \(v\) for \(v\ge5/8\): the factor
\(1+v-\sqrt v\) is increasing for \(v>1/4\), and
\[
 \frac{\dd}{\dd v}\log\frac{(1+v)^{N-1}}{\sqrt{2+v}}
 =\frac{N-1}{1+v}-\frac1{2(2+v)}>0.
\]
After division by \(N(N+1)\), its ratio under \(N\mapsto N+2\) is at least
\[
 \left(\frac{13}{8}\right)^2
 \frac{N(N+1)(N+4)}{(N+2)^2(N+3)}>1,
\]
because, after clearing denominators, the excess is
\[
       105N^3+397N^2-348N-768>0
       \qquad(N\ge7).
\]
It remains to check \(N=7\), \(v=5/8\).  The estimates
\(\sqrt{10}<19/6\) and \(\sqrt{21}<23/5\) give
\[
 \frac9{\sqrt{21/4}}
 \left(\frac{13}{8}-\sqrt{\frac58}\right)
 \left(\frac{13}{8}\right)^6
 >\frac{75}{23}\left(\frac{13}{8}\right)^6
 >56,
\]
the final excess being \(24369203/6029312\).  Thus
\eqref{eq:large-v-target} holds.
\end{proof}

\subsection{Below the cycle scale: the growth lemma}

It remains, except for \(N=7\), to handle \(\zeta\le N\).  The compensation
lemma reduces this to one variable.

\begin{lemma}[One-variable growth]\label{lem:J-growth}
Let \(N\ge9\) be odd.  If \(0\le x<1\) and
\begin{equation}\label{eq:J-domain}
       1-x\ge\frac{2x^2}{N},
\end{equation}
then
\begin{equation}\label{eq:J-ineq}
 J_N(x):=
 \frac{N+2}{N}(1-x)^2(1+x^3)
 \left(\frac{N(1+x^2)}{N-x^2}\right)^{3N/4-1}
 \ge1.
\end{equation}
\end{lemma}

\begin{proof}
Put \(A=3N/4-1\).

\smallskip
\noindent\emph{Step 1: \(0\le x\le2/5\).}
Taking logarithms in \eqref{eq:J-ineq}, and using
\begin{align*}
 \log(1+2/N)&\ge\frac2N-\frac2{N^2},\\
 \log(1+x^2)&\ge x^2-\frac{x^4}{2},\\
 -\log(1-x^2/N)&\ge\frac{x^2}{N},
\end{align*}
together with
\begin{equation}\label{eq:log-shape-bound}
 2\log(1-x)+\log(1+x^3)
 \ge-2x-x^2-\frac{x^4}{1-x}-\frac{x^6}{2},
\end{equation}
we obtain
\begin{equation}\label{eq:L-N-def}
\begin{aligned}
 \log J_N(x)\ge\mathcal L_N(x):={}&
 \frac2N-\frac2{N^2}-2x\\
 &+\left[A\left(1+\frac1N\right)-1\right]x^2
 -\frac A2x^4-\frac{x^4}{1-x}-\frac{x^6}{2}.
\end{aligned}
\end{equation}
For completeness, \eqref{eq:log-shape-bound} follows by writing
\[
 2\log(1-x)
 =-2x-x^2-\frac23x^3-2\sum_{j\ge4}\frac{x^j}{j}
 \ge-2x-x^2-\frac23x^3-\frac{x^4}{1-x}
\]
and using \(\log(1+x^3)\ge x^3-x^6/2\).

On \(0\le x\le1/4\),
\[
 x^4\le\frac{x^2}{16},
 \qquad
 \frac{x^4}{1-x}\le\frac{x^2}{12},
 \qquad
 \frac{x^6}{2}\le\frac{x^2}{512}.
\]
Thus
\[
       \mathcal L_N(x)\ge c_N-2x+B_Nx^2,
\]
where
\[
 c_N=\frac2N-\frac2{N^2},
 \qquad
 B_N=\frac{1116N^2-2003N-1536}{1536N}.
\]
For \(N\ge9\), \(B_N>0\), and
\begin{equation}\label{eq:discriminant-check}
 B_Nc_N-1
 =\frac{348N^3-3119N^2+467N+1536}{768N^3}>0.
\end{equation}
The numerator equals \(6792\) at \(N=9\); its derivative is \(28889\)
there and is increasing afterwards.  Hence the quadratic has positive global
minimum, and \(\mathcal L_N(x)>0\) on \([0,1/4]\).

On \(1/4\le x\le2/5\), differentiation in the real parameter \(N\) gives
\[
 \frac{\partial\mathcal L_N}{\partial N}
 \ge\frac{93}{2048}-\frac2{81}>0.
\]
It is therefore enough to take \(N=9\).  Direct differentiation gives
\begin{equation}\label{eq:L9-data}
 \mathcal L_9(1/4)=\frac{11795}{663552},
 \qquad
 \mathcal L_9'(1/4)=\frac{1295}{3072},
 \qquad
 \mathcal L_9''(2/5)=\frac{49}{3375}.
\end{equation}
Moreover,
\[
 \mathcal L_9'''(x)=-\frac{3xR(x)}{(1-x)^4},
\]
where
\[
 R(x)=20x^6-80x^5+143x^4-174x^3+166x^2-104x+31.
\]
For \(0\le x\le2/5\),
\[
 R(x)\ge31-104x+\frac{2282}{25}x^2
 \ge\frac{2503}{625}>0.
\]
Indeed, the difference from the first quadratic is
\[
 \frac{x^2}{25}
 (500x^4-2000x^3+3575x^2-4350x+1868).
\]
For \(x\le2/5\), the polynomial in parentheses is at least
\[
 500x^4+(3575-800)x^2+(1868-1740)>0,
\]
because \(-2000x^3\ge-800x^2\) and \(-4350x\ge-1740\).
The quadratic \(31-104x+(2282/25)x^2\) is decreasing on this interval and
has value \(2503/625\) at \(x=2/5\).  Thus
\(\mathcal L_9'''<0\).  Since \(\mathcal L_9''(2/5)>0\), we have
\(\mathcal L_9''>0\) throughout the interval; the first two inequalities in
\eqref{eq:L9-data} then imply \(\mathcal L_9>0\).

\smallskip
\noindent\emph{Step 2: \(2/5\le x\le\beta_N\).}
Here \(\beta_N\) is the positive root of
\[
       1-\beta_N=\frac{2\beta_N^2}{N};
\]
condition \eqref{eq:J-domain} is exactly \(x\le\beta_N\).  Since
\(N/(N-x^2)>1\), it is enough to prove
\begin{equation}\label{eq:K-N-def}
 K_N(x):=
 \frac{N+2}{N}(1-x)^2(1+x^3)(1+x^2)^{3N/4-1}\ge1.
\end{equation}
The logarithmic derivative of \(K_N\) has the opposite sign to
\begin{align}
 P_N(x)={}&3Nx^5-3Nx^4+3Nx^2-3Nx
          +6x^5-2x^4\\
         &+10x^3-6x^2+4x+4.
 \label{eq:P-N}
\end{align}
At \(x=2/5\),
\[
       P_N(2/5)=-\frac{2(1197N-8266)}{3125}<0,
\]
while
\[
       P_N''(2/5)=\frac{6(17N+66)}{25}>0
\]
and
\[
 P_N'''(x)=12\left[N(15x^2-6x)+30x^2-4x+5\right]>0
 \qquad(x\ge2/5).
\]
Thus \(P_N\) is convex on the interval and, starting negative, can cross
zero at most once.  Therefore \(K_N\) has no interior minimum: it is enough
to check \(x=2/5\) and \(x=\beta_N\).

At \(x=2/5\),
\[
       K_9(2/5)=\frac{1463}{3125}
       \left(\frac{29}{25}\right)^{23/4}.
\]
Taylor's formula with positive third derivative gives
\[
 \left(1+\frac4{25}\right)^{23/4}
 \ge1+\frac{23}{4}\frac4{25}
 +\frac12\frac{23}{4}\frac{19}{4}
       \left(\frac4{25}\right)^2
 =\frac{2837}{1250}.
\]
Hence
\[
       K_9(2/5)\ge\frac{4150531}{3906250}>1.
\]
Furthermore,
\[
 \frac{K_{N+2}(2/5)}{K_N(2/5)}
 =\frac{N(N+4)}{(N+2)^2}
  \left(\frac{29}{25}\right)^{3/2}
 \ge\frac{117}{121}\frac{31}{25}>1,
\]
where Bernoulli's inequality was used in the last step.  Thus the first
endpoint is settled for every odd \(N\ge9\).

For the second endpoint, use
\((1-\beta_N)^2=4\beta_N^4/N^2\).  The function
\(h_N(x)=1-x-2x^2/N\) is strictly decreasing, and \(h_N(\beta_N)=0\).
The exact evaluations
\[
 h_9(21/25)=\frac2{625},\qquad
 h_{11}(6/7)=\frac5{539},\qquad
 h_{13}(7/8)=\frac3{416}
\]
show that \(\beta_9>21/25\), \(\beta_{11}>6/7\), and
\(\beta_N\ge7/8\) for every \(N\ge13\).  When \(N=9\), hence
\[
\begin{aligned}
 K_9(\beta_9)
 &>\frac{11}{9}\frac4{81}
    \frac{497}{1000}\frac{159}{100}
    \left(\frac{17}{10}\right)^5\frac{37}{25}\\
 &=\frac{15221984467459}{15187500000000}>1.
\end{aligned}
\]
Here we used
\[
 \left(\frac{21}{25}\right)^4>\frac{497}{1000},
 \quad
 1+\left(\frac{21}{25}\right)^3>\frac{159}{100},
 \quad
 1+\left(\frac{21}{25}\right)^2>\frac{17}{10},
\]
and
\[
       \left(\frac{17}{10}\right)^{3/4}>\frac{37}{25}.
\]
The last inequality follows after raising both sides to the fourth power.

When \(N=11\), \(\beta_{11}>6/7\), and the coarser estimates
\(\beta_{11}^4>1/2\), \(1+\beta_{11}^3>3/2\), and
\(1+\beta_{11}^2>17/10\) give
\[
 K_{11}(\beta_{11})
 >\frac{13}{11}\frac4{121}\frac12\frac32
   \left(\frac{17}{10}\right)^7
 =\frac{16003208247}{13310000000}>1.
\]

Finally, for \(N\ge13\), \(\beta_N\ge7/8\).  Dropping the factors
\((N+2)/N\) and \(1+\beta_N^3\),
\[
 K_N(\beta_N)
 \ge\frac{4(7/8)^4}{N^2}
    \left(\frac{113}{64}\right)^{3N/4-1}.
\]
At \(N=13\), this is larger than
\[
       \frac{4(7/8)^4}{13^2}\left(\frac74\right)^8
       =\frac{13841287201}{11341398016}>1.
\]
The ratio of these lower bounds under \(N\mapsto N+2\) is at least
\[
       \left(\frac{13}{15}\right)^2
       \left(\frac74\right)^{3/2}>1.
\]
This completes all endpoint checks and the proof of
\eqref{eq:J-ineq}.
\end{proof}

\begin{lemma}[Below the cycle scale]\label{lem:linear-low-zeta}
If \(\zeta\le N\) and \(N\ge9\) is odd, then
\eqref{eq:linear-core} holds.
\end{lemma}

\begin{proof}
By \Cref{lem:compensation},
\begin{align}
 &(N+2)c_\xi(1-\ell)(\sigma-\ell)\sigma^{N-1}\notag\\
 &\quad\ge
 (N+2)(1-\ell)^2(1+\ell^3)
 \sigma^{N-1-\zeta/4}.
 \label{eq:linear-comp-expanded}
\end{align}
Indeed, \((1+\ell)(\sigma-\ell)\ge1+\ell^3\) because
\(\sigma\ge1+\ell^2\).

Since \(\zeta=\ell^2/u\le N\),
\[
       u\ge\frac{\ell^2}{N},
       \qquad
       a\le\frac{N-\ell^2}{N}.
\]
Thus
\begin{equation}\label{eq:sigma-lower-J}
       \sigma=\frac{1+\ell^2}{a}
       \ge\frac{N(1+\ell^2)}{N-\ell^2}.
\end{equation}
Also \(N-1-\zeta/4\ge3N/4-1\).  Finally,
\(\ell<2a-1\) from the proof of \eqref{eq:T-bound-two}, so
\[
       1-\ell>2u=\frac{2\ell^2}{\zeta}
       \ge\frac{2\ell^2}{N}.
\]
Applying \Cref{lem:J-growth} with \(x=\ell\) to
\eqref{eq:linear-comp-expanded} shows that its right side is at least
\(N\).  Since \(T_N<N\), \eqref{eq:linear-core} follows.
\end{proof}

\subsection{The exceptional cycle scale \texorpdfstring{\(N=7\), i.e. \(m=9\)}{N=7, i.e. m=9}}

The preceding lemmas already cover \(N=7\) whenever either
\(\zeta\ge7\) or \(v\ge5/8\).  It remains to prove
\eqref{eq:linear-core} under
\begin{equation}\label{eq:N7-corner}
       N=7,
       \qquad \zeta\le7,
       \qquad v\le\frac58.
\end{equation}

\begin{lemma}[The small-\(v\) part of \(N=7\)]\label{lem:N7-small-v}
Under \eqref{eq:N7-corner}, if \(0<v\le1/4\), then
\eqref{eq:linear-core} holds.
\end{lemma}

\begin{proof}
Put \(s=\sqrt v\).  Since
\[
       \frac{\ell^2}{v}=\frac{\zeta}{\zeta+1+v}
       \le\frac7{8+v}<\frac78<\left(\frac{15}{16}\right)^2,
\]
we have
\begin{equation}\label{eq:ell-small-v}
       \ell<\frac{15}{16}s.
\end{equation}

We next prove
\begin{equation}\label{eq:xi-small-v}
       \xi\ge\frac1{4v}.
\end{equation}
Let \(t=\zeta+v\).  Formula \eqref{eq:xi-zeta-v} becomes
\[
       4v\xi=
       \frac{16(t+1)^2}{t^2(t(v+2)+2)}.
\]
The reciprocal factor
\[
       G(t,v)=\frac{t^2(t(v+2)+2)}{(t+1)^2}
\]
is increasing in both variables, since
\[
 \frac{\partial G}{\partial t}
 =\frac{t(t^2v+2t^2+3tv+6t+4)}{(t+1)^3}>0,
 \qquad
 \frac{\partial G}{\partial v}=\frac{t^3}{(t+1)^2}>0.
\]
Since
\(t\le7+1/4=29/4\), its maximum is
\[
       \frac{246413}{17424}<16.
\]
This proves \eqref{eq:xi-small-v}.

The function \((1+4\xi)/(2+4\xi)\) is increasing in \(\xi\), so
\eqref{eq:xi-small-v} gives
\[
       \varphi\ge\frac{1+v}{1+2v}.
\]
Also \(\alpha>1/(2+v)\), hence
\[
 c_\xi\sigma^6
 \ge\frac{(1+v)^7}{(1+2v)\sqrt{1+v/2}}.
\]
We shall use the elementary lower bound
\begin{equation}\label{eq:cxi-small-v-poly}
 \frac{(1+v)^7}{(1+2v)\sqrt{1+v/2}}
 \ge\left(1-\frac{7v}{8}\right)(1+5v+11v^2).
\end{equation}
Both sides are positive.  If \(L(v)\) and \(R(v)\) denote respectively
the left and right sides of \eqref{eq:cxi-small-v-poly}, then exact
expansion gives
\begin{equation}\label{eq:small-v-square-diff}
\begin{aligned}
 &128(1+2v)^2(1+v/2)\bigl(L(v)^2-R(v)^2\bigr)\\
 &\quad=v\bigl(128v^{13}+1792v^{12}+11648v^{11}+46592v^{10}
 +128128v^9\\
 &\qquad\quad+232540v^8+345884v^7+492971v^6+464196v^5\\
 &\qquad\quad+258097v^4+86924v^3+18035v^2+2254v+160\bigr),
\end{aligned}
\end{equation}
which is nonnegative.

For the defect, use \eqref{eq:zeta-simple-lower}.  From
\eqref{eq:T-bound-one},
\[
\begin{aligned}
 T_7
 &\le7+7v-(1-v/2)(1-v^{7/2})\\
 &\le6+\frac{15}{2}v+\frac12v^3,
\end{aligned}
\]
because \(\sqrt v\le1/2\).  On the other hand,
\eqref{eq:ell-small-v} and \eqref{eq:cxi-small-v-poly} show that the right
side of \eqref{eq:linear-core} is at least
\[
 9\left(1-\frac78s^2\right)
   \left(1-\frac{15}{16}s\right)
   \left(1+s^2-\frac{15}{16}s\right)
   (1+5s^2+11s^4).
\]
After subtracting \(6+15s^2/2+s^6/2\), we obtain the polynomial
\begin{equation}\label{eq:P9-small}
\begin{aligned}
 P_9(s)={}&\frac{10395}{128}s^9-\frac{333333}{2048}s^8
 +\frac{13635}{128}s^7+\frac{51005}{2048}s^6\\
 &-\frac{18765}{128}s^5+\frac{264969}{2048}s^4
 -\frac{4995}{64}s^3+\frac{11913}{256}s^2\\
 &-\frac{135}{8}s+3.
\end{aligned}
\end{equation}
By \Cref{lem:bernstein-P9}, \(P_9(s)>0\) on \([0,1/2]\).  This proves
\eqref{eq:linear-core}.
\end{proof}

\begin{lemma}[The middle-\(v\) part of \(N=7\)]\label{lem:N7-middle-v}
Under \eqref{eq:N7-corner}, if \(1/4\le v\le5/8\), then
\eqref{eq:linear-core} holds.
\end{lemma}

\begin{proof}
We first show \(\xi>1/3\).  With \(t=\zeta+v\),
\[
       \xi=\frac4v\frac{(t+1)^2}{t^2(t(v+2)+2)}.
\]
The quantity \(vG(t,v)\), with \(G\) as in the proof of
\Cref{lem:N7-small-v}, is increasing in both variables.  Indeed,
\[
 \frac{\partial(vG)}{\partial t}=v\frac{\partial G}{\partial t}>0,
 \qquad
 \frac{\partial(vG)}{\partial v}
 =\frac{2t^2(tv+t+1)}{(t+1)^2}>0.
\]
At
\((t,v)=(61/8,5/8)\), it equals
\[
       \frac{26214445}{2437632}<12.
\]
Thus \(\xi>1/3\), and hence
\begin{equation}\label{eq:varphi-7/10}
       \varphi\ge\frac7{10}.
\end{equation}

Put
\begin{equation}\label{eq:y-middle}
       y^2=\frac{7v}{8+v}.
\end{equation}
Since \(\zeta\le7\), we have \(\ell\le y\).  The first half of the
compensation proof, \eqref{eq:sqrt-compensation}, gives
\[
       \sqrt{2\alpha}\ge1-\ell^2\ge1-y^2.
\]
Also, \(v^{7/2}\le v/2\) on \([1/4,5/8]\), because
\(v^{5/2}\le(5/8)^{5/2}<1/2\).  Together with
\eqref{eq:zeta-simple-lower}, \eqref{eq:T-bound-one} yields
\begin{equation}\label{eq:T7-middle}
       T_7\le6+8v.
\end{equation}
By \eqref{eq:varphi-7/10}, it is enough to prove
\begin{equation}\label{eq:middle-target}
 \frac{63}{10}(1-y^2)(1-y)(1+v-y)(1+v)^6
 \ge6+8v.
\end{equation}

We verify that the ratio of the two sides is increasing.  Solving
\eqref{eq:y-middle} gives
\[
       v=v(y)=\frac{8y^2}{7-y^2}.
\]
For \(1/4\le v\le5/8\),
\[
       \frac9{20}<y<\frac57.
\]
Define
\[
 \mathcal R(y)=
 \frac{(1-y^2)(1-y)(1+v(y)-y)(1+v(y))^6}
      {6+8v(y)}.
\]
A direct logarithmic differentiation gives
\begin{equation}\label{eq:R-log-derivative}
 \frac{\mathcal R'(y)}{\mathcal R(y)}
 =\frac{2Q_{10}(y)}
 {(y-1)(y+1)(y^2-7)(y^2+1)(29y^2+21)
  (y^3+7y^2-7y+7)},
\end{equation}
where
\begin{align}
 Q_{10}(y)={}&58y^{10}+319y^9-1793y^8-11280y^7+9788y^6\\
 &-11354y^5+350y^4+5880y^3-6174y^2+5635y-1029.
 \label{eq:Q10}
\end{align}
The denominator in \eqref{eq:R-log-derivative} is positive on
\([9/20,5/7]\), and \Cref{lem:bernstein-Q10} gives
\(Q_{10}(y)>0\) there.  Thus \(\mathcal R\) is increasing.

At \(v=1/4\), one has \(y_0=\sqrt{7/33}<6/13\).  Since the factors involving
\(y\) decrease with \(y\),
\begin{align*}
 &\frac{63}{10}(1-y_0^2)(1-y_0)(5/4-y_0)(5/4)^6\\
 &\quad>
 \frac{63}{10}\frac{26}{33}\frac7{13}\frac{41}{52}
 \left(\frac54\right)^6
 =\frac{18834375}{2342912}>8.
\end{align*}
This proves \eqref{eq:middle-target} at the left endpoint.  The monotonicity
of \(\mathcal R\) proves it throughout the interval.
\end{proof}

\begin{proposition}[Linear branch]\label{prop:linear-branch}
For every odd \(N\ge7\), assumption \eqref{eq:linear-branch} implies
\[
       R_m\le C_m\psi(\xi,\rho).
\]
\end{proposition}

\begin{proof}
For \(N\ge9\), if \(\zeta\le N\), use
\Cref{lem:linear-low-zeta}.  If \(\zeta\ge N\), use
\Cref{lem:linear-high-zeta} when \(v\le5/8\), and
\Cref{lem:linear-large-v} when \(v\ge5/8\).

For \(N=7\), the same two broad lemmas cover the cases
\(\zeta\ge7\) or \(v\ge5/8\).  The remaining corner
\eqref{eq:N7-corner} is covered by
\Cref{lem:N7-small-v,lem:N7-middle-v}.
In every case \eqref{eq:linear-core} holds.  Combining it with
\eqref{eq:linear-payment-start}, \eqref{eq:KL-first-term}, and
\eqref{eq:F-T} proves the proposition.
\end{proof}


\section{Assembly of the complete proof}\label{sec:assembly}

\begin{proof}[Proof of \Cref{thm:main}]
Put \(q=1-p\).

If \(p\le1/2\), the right side of \eqref{eq:main-goal} is nonpositive, so the
claim follows from \(t(C_m,W)\ge0\).  Assume henceforth \(p>1/2\).

The cases \(m=3,5,7\) follow at every density from
\Cref{prop:C3,prop:C5,prop:C7}.  Let \(m\ge9\) be odd.

If \(p\ge2/3\), or equivalently \(q\le1/3\), the result is
\Cref{thm:dense-region}.  It remains to consider
\(1/2<p<2/3\), equivalently \(1/3<q<1/2\).
By \Cref{lem:step-reduction}, it is enough to treat a step graphon.  If the
compression \(A\) has no eigenvalue above \(q\),
\Cref{lem:no-frontier} proves the theorem.  Otherwise there is a unique
frontier eigenvalue \(\alpha\), and
\Cref{prop:master-defect,thm:huber-elim} reduce the problem to the scalar
target \eqref{eq:scalar-target}.  If \(2\rho\xi\le1\),
\Cref{prop:quadratic-branch} proves that target.  If \(2\rho\xi>1\),
\Cref{prop:linear-branch} proves it.  The two alternatives are exhaustive.
The step-graphon approximation returns the result to arbitrary graphons.
\end{proof}

\begin{remark}[What is finite and what is uniform]
The dense-region argument is uniform in every odd \(m\ge3\).  The
intermediate-region argument is uniform in every odd \(m\ge9\), except for a
small corner at \(m=9\), where two fixed univariate polynomials are certified
positive in \Cref{app:bernstein}.  There is no length-by-length computation
for \(m=11,13\), and no adaptive interval subdivision anywhere in the proof.
\end{remark}

\begin{remark}[Equality and stability]
The manuscript is organized to prove the inequality rather than classify all
equality cases.  A stability analysis would require tracking near-equality
in the cycle-to-path deletion, the gamma moment inequality, the
Hilbert--Schmidt budget, the forced-variance inequality, and the active Huber
witness.  The reductions here identify those loci but do not pursue the
classification.
\end{remark}

\section{Transferable mechanisms and the orthogonal-polynomial viewpoint}
\label{sec:transferable}

This section is not needed for logical completeness.  Its purpose is to
extract methods that may be useful in other spectral extremal problems.

\subsection{One-step stripping: the Schur complement as a Jacobi recursion}

Let \(K\) be a finite-dimensional self-adjoint operator, let \(e\) be a unit
vector, and decompose
\[
 K=\begin{pmatrix}a&h^*\\h&B\end{pmatrix}
 \quad\text{on}\quad \R e\oplus e^\perp.
\]
The same block inversion used in \eqref{eq:dense-path-gf} gives
\begin{equation}\label{eq:general-stripping}
 M_{K,e}(z):=\ip e{(I-zK)^{-1}e}
 =\frac1{1-az-z^2\ip h{(I-zB)^{-1}h}}.
\end{equation}
If \(h\ne0\), let \(\nu\) be the probability spectral measure of \(B\) at
\(h/\norm h\).  Then
\begin{equation}\label{eq:general-stripping-normalized}
       M_{K,e}(z)^{-1}
       =1-az-\norm h^2z^2M_\nu(z).
\end{equation}
This is exactly the first step of the Jacobi continued fraction.  Iterating
on the cyclic Krylov space generated by \(e\) produces a tridiagonal Jacobi
matrix and its full J-fraction.  In the centered, variance-one normalization,
\eqref{eq:general-stripping-normalized} becomes
\[
       1-M_\mu(z)^{-1}=z^2M_\nu(z),
\]
the classical once-stripped moment relation emphasized by
Anshelevich~\cite{Anshelevich2009}; coefficient stripping for scalar and
matrix Jacobi operators is surveyed by Damanik--Pushnitski--Simon
\cite{DamanikPushnitskiSimon2008}.

In our graphon proof, \(e=\one\), \(a=q\), \(h=g\), and \(B=A\).  Thus the
complement degree fluctuation \(g\) is the first off-diagonal Jacobi
coefficient, while \(A\) is the once-stripped tail.  The proof does not need
to tridiagonalize the tail; it retains its spectral measure \(\mu\) directly.
This is useful in infinite-dimensional problems where the first Jacobi step
is rigid but the tail is not.

\subsection{First returns, irreducible excursions, and logarithmic cycles}

Equation \eqref{eq:dense-path-gf} has the first-return form
\[
       M(z)=\frac1{1-H(z)},
       \qquad H(z)=qz+z^2R_+(z).
\]
Combinatorially, a closed walk from the distinguished state is a sequence of
irreducible returns: either a level step of weight \(q\), or a departure into
the orthogonal tail followed by a return, encoded by \(z^2R_+\).  This is
the same structural decomposition that underlies Flajolet's continued
fraction for weighted Motzkin paths~\cite{Flajolet1980} and Lehner's use of
Jacobi matrices, lattice paths, and irreducible paths in cumulant
expansions~\cite{Lehner2003}.

The logarithms \(-\log(1-Y_\pm)=\sum_{r\ge1}Y_\pm^r/r\) perform a related
but cyclic assembly: \(r\) excursions are concatenated with the cyclic
weight \(1/r\).  The two-sided identity then applies a parity projection:
for odd coefficients, the stripped-tail terms from \(A\) and \(-A\) cancel.
A general design principle suggested by this proof is:
\begin{quote}
Separate a distinguished state, encode returns by a Schur complement, and
combine a model with a signed or complemented companion so that the
uncontrolled tail moments cancel in the desired parity class.
\end{quote}
This ``strip, assemble, cancel'' pattern is broader than graphons.

\subsection{Dirichlet diagonalization as probabilistic polarization}

The calculation in \Cref{lem:dense-dirichlet} is an instance of the following
general lemma.

\begin{proposition}[Dirichlet polarization]\label{prop:general-dirichlet}
Fix \(r\ge1\).  Suppose
\[
       F(\lambda_1,\ldots,\lambda_r)
       =\sum_{j=0}^d c_jh_j(\lambda_1,\ldots,\lambda_r),
\]
and define
\[
       \widetilde F(\ell)
       =\sum_{j=0}^d c_j\binom{j+r-1}{r-1}\ell^j.
\]
If \(\Theta\sim\operatorname{Dirichlet}(1,\ldots,1)\), then
\[
       F(\lambda_1,\ldots,\lambda_r)
       =\E\widetilde F(\Theta\cdot\lambda).
\]
Consequently, for every interval \(I\subset\R\),
\[
       \min_{\lambda\in I^r}F(\lambda)
       =\min_{\ell\in I}\widetilde F(\ell).
\]
\end{proposition}

\begin{proof}
This is exactly \eqref{eq:dense-dirichlet-moment}, followed by linearity.
The minimum statement follows because \(\Theta\cdot\lambda\in I\), and the
reverse inequality is obtained on the diagonal.
\end{proof}

This lemma may be viewed as a positive, probabilistic version of
polarization: the complete homogeneous basis is precisely the basis for
which restriction to the diagonal can be inverted by averaging over a
simplex.  The law of \(\Theta\cdot\lambda\) is a Dirichlet average and, in
one dimension, is closely related to B-splines and the Hermite--Genocchi
formula for divided differences; see de~Boor~\cite{deBoor2005}.  The key
extremal consequence is unusually strong: positivity on a whole spectral
cube is equivalent to positivity on the diagonal whenever the polynomial
has this basis structure.

\subsection{Beta--gamma algebra and Laplace positivity transfer}

The passage from \eqref{eq:dense-beta-integral} to
\eqref{eq:dense-gamma-smoothing} is an instance of a general positive
transform.  For suitable \(f\), integers \(m>r\ge1\), and \(\ell>0\),
\begin{equation}\label{eq:general-laplace-transfer}
 \int_0^\infty\frac{s^{r-1}}{(\ell+s)^m}f(s)\dd s
 =\frac{\Gamma(r)}{\Gamma(m)}
  \int_0^\infty t^{m-r-1}e^{-\ell t}
       \E f(Y/t)\dd t,
 \qquad Y\sim\GammaLaw(r,1).
\end{equation}
The proof is just the Laplace representation of \((\ell+s)^{-m}\), followed
by \(y=ts\).  Its value is conceptual: a parameter-dependent rational
kernel becomes a positive mixture of expectations under a fixed
probability family.  Therefore any pointwise-in-\(t\) expectation inequality
for \(f(Y/t)\) transfers automatically to every \(\ell\).

In the present proof, the repeated variables in the complete homogeneous
polynomials first produce a beta distribution; the projective coordinate
\(s=\ell(1-x)/x\) turns the beta density into a beta-prime kernel; and the
Laplace transform reveals a gamma variable.  This beta--gamma route is a
useful search heuristic whenever repeated spectral parameters and rational
powers occur together.

\subsection{The gamma Stein identity and a hidden Laguerre equation}

Equation \eqref{eq:gamma-Stein} is the standard gamma Stein identity; the
operator
\[
       \mathcal A_rf(y)=yf'(y)+(r-y)f(y)
\]
characterizes the \(\GammaLaw(r,1)\) law and is widely used in Stein's
method; see, for example, Gaunt~\cite{Gaunt2016}.  Here it is used not for
distributional approximation but as a ladder relation for shifted moments.

There is an exact orthogonal-polynomial shadow.  For an integer \(k\ge0\),
put
\[
       F_k(b)=\E(Y-b)^k.
\]
Expansion of the moment, or the Rodrigues formula, gives
\begin{equation}\label{eq:Laguerre-identity}
       F_k(b)=(-1)^k k!\,L_k^{(-r-k)}(-b),
\end{equation}
where \(L_k^{(\alpha)}\) is the generalized Laguerre polynomial.  For
\(k=2j\), the Laguerre differential equation is exactly
\eqref{eq:gamma-ODE}.  The parameter \(-r-2j\) lies outside the classical
orthogonality range \(\alpha>-1\), so we are using the algebraic differential
and ladder structure of Laguerre polynomials, not a positive Laguerre
orthogonality measure.  Standard background is in Szeg\H{o}
\cite{Szego1975}.

The function \(G\) in \eqref{eq:gamma-G} is a weighted amplitude whose
critical points factor explicitly.  Proving that its left maximum dominates
its right maximum and then using \eqref{eq:gamma-ibp-max} is analogous in
spirit to Sonin--Sturm arguments for comparing successive extrema of
solutions to second-order equations.  The important transferable move is:
construct a first-order multiplier so that the desired linear combination
of a polynomial solution and its derivative becomes an integral of a
factored derivative; positivity is then reduced to a global maximum
comparison.

This explains the connection with the orthogonal-polynomial techniques
suggested by the two cited papers.  Anshelevich's once-stripped moment
identity is literally our first Schur step, while Lehner's Jacobi-matrix and
Motzkin-path picture explains the return decomposition of the moment series.
The later gamma lemma lands in a different classical family---Laguerre---but
uses the same general philosophy: a recurrence or differential equation is
more effective than coefficientwise estimation of high moments.

\subsection{The one-dangerous-mode principle and Huber duality}

The intermediate region exhibits a complementary mechanism.  The square
budget says that two eigenvalues above \(q\) would cost more than \(pq\) when
\(q>1/3\), so there is at most one dangerous mode.  Entrywise constraints on
the corresponding eigenfunction force a choice: either \(g\) couples
substantially to that mode, or a quantitatively controlled part of \(g\)
must lie in the safe subspace.  Retaining both alternatives leads to
\[
       \rho v^2+(\xi-v+v^2)_+,
\]
and minimization over the shape variable \(v\) produces the Huber envelope
\(\psi(\xi,\rho)\).  Its dual variable \(\lambda\in[0,1]\) is an active-set
multiplier.  The witnesses \(\lambda=2\rho\xi\) and \(\lambda=1\) are the
quadratic and saturated linear regimes.

This suggests a general recipe for constrained spectral inequalities:
\begin{enumerate}[leftmargin=2.2em]
\item use a trace or energy budget to isolate a finite set of dangerous
      modes;
\item derive geometric constraints that force compensation either in direct
      coupling or in the safe complement;
\item keep the channels simultaneously rather than estimating them
      separately;
\item eliminate the shape by convex duality, and choose dual witnesses
      adapted to small and large activation.
\end{enumerate}
The method should extend to several dangerous modes, where the scalar Huber
envelope is replaced by a low-dimensional convex program.

\subsection{Why \texorpdfstring{\(q=1/3\)}{q=1/3} is the natural interface}

The same threshold appears for independent structural reasons.

On the dense side, centering at \(1/2\) gives
\(\delta=1/2-q\).  The gamma moment lemma pays for the odd term in
\eqref{eq:dense-rho-centered} precisely when
\[
       2\delta\ge\frac13,
       \qquad\text{i.e.}\qquad q\le\frac13.
\]
On the intermediate side, two eigenvalues larger than \(q\) would consume
more than \(2q^2\) of the square budget, and uniqueness is guaranteed
precisely when
\[
       2q^2>pq=q(1-q),
       \qquad\text{i.e.}\qquad q>\frac13.
\]
Thus the two mechanisms meet without overlap or gap.  This matching strongly
suggests that \(q=1/3\) is an intrinsic phase boundary of the proof problem,
not an artifact of constants chosen later in the argument.
\appendix

\section{The two Bernstein certificates}\label{app:bernstein}

For completeness, this appendix records the only two fixed polynomial
positivity checks used in the proof.  The Bernstein-basis principle used
here is reviewed in Farouki~\cite{Farouki2012}.  We include the coefficients themselves,
not merely a reference to a computer calculation.

For an interval \([a,b]\), put
\[
        t=\frac{x-a}{b-a},
        \qquad
        B^{[a,b]}_{i,n}(x)=\binom ni t^i(1-t)^{n-i}
        \quad(0\le i\le n).
\]
The functions \(B^{[a,b]}_{i,n}\) are nonnegative on \([a,b]\) and sum to
one.  Hence an expansion
\[
        P(x)=\sum_{i=0}^n \beta_i B^{[a,b]}_{i,n}(x)
\]
with all \(\beta_i>0\) proves \(P(x)>0\) throughout the interval.  The
coefficients can be checked directly as follows.  If
\[
        P(a+(b-a)t)=\sum_{j=0}^n c_jt^j,
\]
then
\begin{equation}\label{eq:power-to-bernstein}
        \beta_i=
        \sum_{j=0}^i c_j\frac{\binom ij}{\binom nj}
        \qquad(0\le i\le n).
\end{equation}
Indeed,
\[
        t^j=\frac1{\binom nj}
             \sum_{i=j}^n\binom ij B_{i,n}(t),
\]
and substitution gives \eqref{eq:power-to-bernstein}.

\begin{lemma}[The degree-nine certificate]\label{lem:bernstein-P9}
The polynomial in \eqref{eq:P9-small}, namely
\begin{align*}
 P_9(s)={}&
 \frac{10395}{128}s^9
 -\frac{333333}{2048}s^8
 +\frac{13635}{128}s^7
 +\frac{51005}{2048}s^6\\
 &-\frac{18765}{128}s^5
 +\frac{264969}{2048}s^4
 -\frac{4995}{64}s^3
 +\frac{11913}{256}s^2
 -\frac{135}{8}s+3,
\end{align*}
is positive on \([0,1/2]\).
\end{lemma}

\begin{proof}
Its degree-nine Bernstein coefficients on \([0,1/2]\), in increasing order
of the index, are
\begin{align*}
 &3,
 &&\frac{33}{16},
 &&\frac{17795}{12288},
 &&\frac{29843}{28672},
 &&\frac{361761}{458752},\\[1mm]
 &\frac{918263}{1376256},
 &&\frac{7142309}{11010048},
 &&\frac{1096445}{1572864},
 &&\frac{390495}{524288},
 &&\frac{361511}{524288}.
\end{align*}
Every numerator and denominator is positive.  The Bernstein
partition-of-unity argument therefore proves the claim.  Formula
\eqref{eq:power-to-bernstein} gives a short exact-arithmetic verification of
the displayed list.
\end{proof}

\begin{lemma}[The degree-ten certificate]\label{lem:bernstein-Q10}
The polynomial \(Q_{10}\) in \eqref{eq:Q10} is positive on
\([9/20,5/7]\).
\end{lemma}

\begin{proof}
Its degree-ten Bernstein coefficients on \([9/20,5/7]\), in increasing
order of the index, are
\begin{align*}
 &\frac{3243737479716939}{5120000000000},
 &&\frac{24439395309190297}{35840000000000},
 &&\frac{2256629681468091}{3136000000000},\\[1mm]
 &\frac{653354162075429}{878080000000},
 &&\frac{3033296182962901}{4033680000000},
 &&\frac{625950152734243}{847072800000},\\[1mm]
 &\frac{1151767076423}{1647086000},
 &&\frac{6182915802683}{9882516000},
 &&\frac{7481469276}{14706125},\\[1mm]
 &\frac{67938971678}{201768035},
 &&\frac{26748047904}{282475249}.
\end{align*}
Again every coefficient is positive.  Since the Bernstein basis is
nonnegative and sums to one, \(Q_{10}>0\) on the stated interval.
The conversion formula \eqref{eq:power-to-bernstein} verifies the list using
only rational arithmetic.
\end{proof}

\begin{remark}[Why this is a proof certificate]
A Bernstein-coefficient list is not a numerical sample.  Once the exact
identity with the Bernstein expansion is checked, positivity at every point
of the whole interval follows from the nonnegativity of the basis functions.
In the present proof no subdivision is used: one interval suffices for each
of the two polynomials.
\end{remark}


\begin{thebibliography}{99}

\bibitem{Anshelevich2009}
M.~Anshelevich,
\newblock Appell polynomials and their relatives III: conditionally free theory,
\newblock \emph{Illinois Journal of Mathematics} \textbf{53} (2009), 39--66;
\newblock arXiv:0803.4279.

\bibitem{deBoor2005}
C.~de~Boor,
\newblock Divided differences,
\newblock \emph{Surveys in Approximation Theory} \textbf{1} (2005), 46--69.

\bibitem{DamanikPushnitskiSimon2008}
D.~Damanik, A.~Pushnitski, and B.~Simon,
\newblock The analytic theory of matrix orthogonal polynomials,
\newblock \emph{Surveys in Approximation Theory} \textbf{4} (2008), 1--85.

\bibitem{Farouki2012}
R.~T. Farouki,
\newblock The Bernstein polynomial basis: a centennial retrospective,
\newblock \emph{Computer Aided Geometric Design} \textbf{29} (2012),
379--419.

\bibitem{Flajolet1980}
P.~Flajolet,
\newblock Combinatorial aspects of continued fractions,
\newblock \emph{Discrete Mathematics} \textbf{32} (1980), 125--161.

\bibitem{Gaunt2016}
R.~E. Gaunt,
\newblock Products of normal, beta and gamma random variables: Stein
operators and distributional theory,
\newblock \emph{Brazilian Journal of Probability and Statistics}
\textbf{32} (2018), 437--466;
\newblock arXiv:1507.07696.

\bibitem{Goodman1959}
A.~W. Goodman,
\newblock On sets of acquaintances and strangers at any party,
\newblock \emph{American Mathematical Monthly} \textbf{66} (1959),
778--783.

\bibitem{Huber1964}
P.~J. Huber,
\newblock Robust estimation of a location parameter,
\newblock \emph{Annals of Mathematical Statistics} \textbf{35} (1964),
73--101.

\bibitem{Lehner2003}
F.~Lehner,
\newblock Cumulants, lattice paths, and orthogonal polynomials,
\newblock \emph{Discrete Mathematics} \textbf{270} (2003), 177--191;
\newblock arXiv:math/0110031.

\bibitem{Lovasz2012}
L.~Lov\'asz,
\newblock \emph{Large Networks and Graph Limits},
\newblock American Mathematical Society Colloquium Publications, vol.~60,
American Mathematical Society, Providence, RI, 2012.

\bibitem{PowersReznick2000}
V.~Powers and B.~Reznick,
\newblock Polynomials that are positive on an interval,
\newblock \emph{Transactions of the American Mathematical Society}
\textbf{352} (2000), 4677--4692.

\bibitem{Simon2005}
B.~Simon,
\newblock \emph{Trace Ideals and Their Applications}, second edition,
\newblock Mathematical Surveys and Monographs, vol.~120,
American Mathematical Society, Providence, RI, 2005.

\bibitem{Sion1958}
M.~Sion,
\newblock On general minimax theorems,
\newblock \emph{Pacific Journal of Mathematics} \textbf{8} (1958),
171--176.

\bibitem{Szego1975}
G.~Szeg\H{o},
\newblock \emph{Orthogonal Polynomials}, fourth edition,
\newblock American Mathematical Society Colloquium Publications, vol.~23,
American Mathematical Society, Providence, RI, 1975.

\end{thebibliography}

\end{document}
